% 本文件由 LaTeX 排版 Agent 根据有效 translation.json 自动生成。
% author=Augustine of Hippo; work_id=1302; title=《信德、望德与爱德手册》

\chapter{《信德、望德与爱德手册》}

% structure: p0001 source=h1 target=omitted reason=page-title-already-rendered-by-parent

《手册》（Enchiridion）乃是呈给劳伦斯（Laurentius），以答复他的问题。此人身世不详。一份手抄本称他为执事，另一份则称他为罗马城的公证员。他很可能是一位平信徒。\par

圣奥思定写此书是在热罗尼莫去世之后（主历420年9月30日），因为在第87章中他提到了\Quote{可忆念的}热罗尼莫。\par

作者通常称此书为\WorkTitle{论信德、望德与爱德}，因为他是按照这三个主题来论述的（参见\BibleRef{格前13:13}）。在第一主题下，他遵循宗徒信经的顺序，并驳斥了摩尼教、阿波利纳主义、亚略主义和白拉奇主义等异端，但没有指名道姓。在第二主题下，他对天主经做了简要阐述。第三部分则是关于基督徒爱德的论述。\par

圣奥思定在他的\WorkTitle{更正篇}第二卷第六十三章中这样提到此书：\Quote{我也应那位收信人的请求，写了一本关于\WorkTitle{信德、望德与爱德}的书，以便他手中常有一本我的作品，希腊人称之为\WorkTitle{手册}（Enchiridion）。我以为在那本书中，我相当仔细地讨论了应当如何钦崇天主，而神圣圣经将这种知识定义为人的真智慧。该书以\InnerQuote{我无法表达}开头，等等。}\par

% structure: p0006 source=h2 target=section reason=relative-heading-rank; base=h2

\section{引言}

% structure: p0007 source=h3 target=subsection reason=relative-heading-rank; base=h2

\subsection{第一章 作者渴望劳伦斯获得真智慧}

我亲爱的儿子劳伦斯，我无法表达，看到你在知识上的进步，以及我多么诚恳地希望你成为一个智慧人——不是那种被称为\BibleQuote{智慧人在哪里？经师在哪里？此世的辩士在哪里？天主岂不是使此世的智慧变成愚拙吗？}的人，而是那种被称为\BibleQuote{智慧人的群众是世界的福泽}的人，也是宗徒希望他们成为的那种人：他告诉他们\BibleQuote{我愿你们在善事上智慧，在恶事上单纯}。如今，正如没有人能靠自己存在一样，也没有人能靠自己成为智慧，而只能靠那一位的启迪光照，正如经上所记：\BibleQuote{一切智慧皆来自上主。}\par

% structure: p0009 source=h3 target=subsection reason=relative-heading-rank; base=h2

\subsection{第二章 敬畏上主是人的真智慧}

人的真智慧是虔敬（pietas）。你在约伯圣传中可以看到这一点。因为我们在那里读到智慧本身对人所说的话：\BibleQuote{看，敬畏上主（虔敬），就是智慧。}如果你进一步问此处\Quote{虔敬}指什么，希腊文更明确地称之为\OriginalTerm{对天主的敬拜}{\GreekText{θεοσέβεια}}，即对天主的敬拜。希腊人有时也称虔敬为\OriginalTerm{正确的敬拜}{\GreekText{εὐσέβεια}}，意指正确的敬拜，尽管这当然特别指向对天主的敬拜。但当我们定义人的真智慧在于什么时，最适宜的用语是明确表达敬畏上主的那个词。而你，既然希望我用简短的言语讨论伟大的事情，难道还想得到更简短的表达方式吗？或者，你也许希望这个表达本身得到简要的解释，并希望我用简短的论述阐明敬拜天主的正确方式？\par

% structure: p0011 source=h3 target=subsection reason=relative-heading-rank; base=h2

\subsection{第三章 当以信德、望德与爱德敬拜天主}

现在，如果我回答，应当以信德、望德与爱德来敬拜天主，你会立刻说这个回答太简短，并要我简要说明这三种恩宠的对象，即我们当信什么、当望什么、当爱什么。而当我做完这些，你对你信中所提的一切问题都将得到回答。如果你有信件的副本，你可以轻易地查阅并重读；如果没有，当我唤起你的记忆时，你也不难回忆起来。\par

% structure: p0013 source=h3 target=subsection reason=relative-heading-rank; base=h2

\subsection{第四章 劳伦斯提出的问题}

你说，你希望我为你写一本手册，可以常带在身边，包含对以下问题的回答：人生活的首要目的是什么；面对各种异端，他主要应避免什么；宗教在多大程度上受理性支持；当信仰独存时，理性中有什么不能支持信仰；宗教的起点和终点是什么；全部教义体系的总结是什么；大公信仰确定而恰当的基础是什么。毫无疑问，如果你彻底了解信德、望德与爱德的正确对象，你就会知道所有这些问题的答案。因为这些必须是宗教中首要的，甚至唯一的追求对象。反对这些的人，要么与基督之名完全无关，要么是异端者。这些需要由理性来辩护，而理性必须以身体的感官或心灵的直觉为起点。凡我们没有通过身体感官经验到，也不能通过理智达到的，必定要相信那些见证者的见证，正是他们写成了那恰被称为神圣的圣经；并且他们藉着天主的助佑，得以通过身体的感官或理智的领悟，看见或预见了所说的事。\par

% structure: p0015 source=h3 target=subsection reason=relative-heading-rank; base=h2

\subsection{第五章 对这些问题的简要回答}

再者，当心神被那借着爱德而行动的信德初步所浸润时，它便努力借着生活的纯洁达到目睹，因为心地纯洁完美的人认识那不可言喻的美，其全瞻就是至高幸福。这无疑是对你问题的回答：什么是起点，什么是终点？我们始于信德，而成全于目睹。这也是全部教义的总纲。但大公信仰的稳固而恰当的基础是基督。宗徒说：\BibleQuote{因为没有别的根基，}宗徒说，\BibleQuote{人不能立别的根基，所立的就是耶稣基督。}我们也不可否认这是大公信仰的恰当基础，因为可能有些异端者也与我们共同持有这一点。因为我们若仔细思考关于基督的事，就会发现，在那些自称为基督徒的异端者中，基督仅在名号上存在；在行动和真理上，他并不在他们中间。但要证明这一点会花费太长时间，因为我们需要逐一审视所有过去、现在或将来可能以基督徒名义存在的异端，并证明每一条异端都是如此——这样讨论几乎无休无止，需要许多卷册。\par

% structure: p0017 source=h3 target=subsection reason=relative-heading-rank; base=h2

\subsection{第六章 在本手册书中不宜展开辩论。}

现在你向我索要一本手册，即一本可以握在手中的书，而不是一本可以塞满书架的书。那么，回到那三种恩宠——如我所说，天主应当借以受崇拜——即信德、望德、爱德：说明每种恩宠的真正与恰当对象是容易的。但要保卫这一真道，对抗那些持相反意见者的攻击，就需要更充分更详尽的教导。而获得这种教导的真正方式，不是把一篇短论放在手中，而是心中燃起极大的热忱。\par

% structure: p0019 source=h3 target=subsection reason=relative-heading-rank; base=h2

\subsection{第七章 信经和天主经要求信德、望德、爱德的实践。}

因为你有信经和天主经。还有什么比这更简短易听易读？还有什么比这更容易记忆？当人类因罪恶在沉重的苦难重压下呻吟，迫切需要天主的怜悯时，一位先知预见了天主恩宠的时刻，宣告说：\BibleQuote{凡呼号上主名号的人，必然获救。}因此有天主经。但宗徒为了推崇这同一恩宠，引用了这段先知预言后，立刻加上：\BibleQuote{那么，人还没有信他，怎能呼号他呢？}因此有信经。在这两者中，你已看到三种恩宠得到体现：信德相信，望德和爱德祈祷。但若无信德，后两者不能存在，因此我们也可说信德也祈祷。正如经上所载：\BibleQuote{人还没有信他，怎能呼号他呢？}\par

% structure: p0021 source=h3 target=subsection reason=relative-heading-rank; base=h2

\subsection{第八章 信德与望德的区别，以及信德、望德、爱德的相互依存。}

再者，有哪样事物不是信德的对象却可被希望呢？确实，一样不是望德的对象的事物可以被信。例如，哪个真基督徒不信恶人受罚呢？然而这样的人并不希望它。而那人相信那刑罚正悬临于己，因前景而恐惧退缩，更当称为畏惧而非希望。诗人仔细区分了这两种心态，他说：\Quote{准许恐惧者怀抱望德。}另一位通常高明得多的诗人却误用了这个词，他说：\Quote{若我曾能希望承受如此巨大的悲恸。}有些语法学家以此为例说明用词不当，说道：\Quote{他说了\InnerQuote{\Emph{sperare}}[希望]而非\InnerQuote{\Emph{timere}}[畏惧]。}因此，信德可以以恶为对象，也可以以善为对象；因为善与恶皆可被信，而相信它们的信德并非恶，乃是善。此外，信德涉及过去、现在与将来三者。例如，我们信基督曾死——过往的事件；我们信祂坐在天主右侧——现今的状态；我们信祂将审判生者与死者——将来的事。再者，信德适用于自身与他人的境况。例如，每人相信自己存在有开端，并非永恒，并同样相信他人与万物。再者，我们关于宗教事物许多信念不仅关乎他人，也关乎天使。但望德仅以善为对象，仅以将来为对象，仅以影响怀有此望德者为对象。因此，信德与望德必须加以区分，不仅是用词得当的问题，更因它们本质不同。我们既看不见所信的，也看不见所望的——这一点是信德与望德的共同之处。例如，\BibleBook{}{希伯来书}中信德被定义为（天主教信仰的杰出捍卫者以此为准则）\BibleQuote{是未见之事的明证}。然而，若有人说他相信，即其信德并非基于言语、见证或任何推理，而是基于自身感官的直接证据，那么他并不犯有用词不当之过，足以合理受到批评：\Quote{你看见了，所以你没有相信。}因此，信德的对象并非不能是视觉的对象。但我们更宜按照圣经所教导的，使用\Quote{信}一字，将其应用于未见之事。关于望德，宗徒又说：\BibleQuote{所见到的望德不是望德；因为人见到的，为什么还希望呢？但我们若希望那未见的，就要忍耐等待。}因此，当我们相信善将临，这无非就是希望它。我还能怎样论爱德呢？没有爱德，信德毫无益处；爱德若不在，望德便不能存在。雅各伯宗徒说：\BibleQuote{魔鬼也信，且战栗。}——即他们既无望德也无爱德，却相信我们所爱所望的将要来临，因而恐惧。因此保禄宗徒称赞并认可\BibleQuote{以爱德行事的信德}；这信德无疑不能离开望德而存在。因此，没有望德，爱德便不存在；没有爱德，望德便不存在；没有信德，爱德和望德皆不存在。\par

% structure: p0023 source=h2 target=omitted reason=duplicate-page-title

% structure: p0024 source=h3 target=subsection reason=relative-heading-rank; base=h2

\subsection{第九章 我们应相信什么。关于受造界，基督徒无需知道更多，只需知道造物主的善是万有之因。}

因此，当有人问及我们关于宗教应信什么时，不必像希腊人称为\Emph{physici}的人那样探究事物的本性；也不必担心基督徒不知道元素的力量与数目——天体的运动、秩序与蚀；天穹的形状；动植物、石头、泉水、河流、山岳的种类与本性；年代学与距离；暴风雨来临的征兆；以及哲学家们发现或自以为发现的成千上万其他事物。因为这些人们自身，虽有如此天赋、热情洋溢、闲暇充裕，有些靠人类推测追踪，有些借助历史与经验探究，却并未发现一切；他们自诩的发现往往不过是猜测而非确切知识。基督徒足以相信：所有受造之物，无论天上地下、可见不可见，其唯一原因是造物主——唯一真天主的善；并且除祂自身外，无一物不从祂获得存在；祂就是天主圣三——即圣父、由圣父所生的圣子，以及由同一圣父所发，但同为圣父与圣子之神的圣神。\par

% structure: p0026 source=h3 target=subsection reason=relative-heading-rank; base=h2

\subsection{第十章 至善的造物主所造万物皆善。}

由如此至上、同等、不变的善的天主圣三，万物被造；这些事物并非至上、同等、不变的善，但它们确是善，即使单独来看也是如此。然而，合为一体，它们非常好，因为它们的\Emph{整体}构成了宇宙及其奇妙秩序与美丽。\par

% structure: p0028 source=h3 target=subsection reason=relative-heading-rank; base=h2

\subsection{第十一章 宇宙中所谓恶不过是善的缺失。}

在宇宙中，即使被称为恶的事物，当它被规范并置于其适当位置时，也只会增强我们对善的赞叹；因为我们通过将善与恶相比较，更加享受和珍视善。因为全能的天主，就连外邦人也承认，祂对万物拥有至高无上的权能，祂本身是至善的，若非祂如此全能且良善，以至于能由恶中引出善来，祂绝不会允许在祂的受造物中存在任何恶的事物。因为我们所谓的恶，难道不就是善的缺失吗？在动物的肉体中，疾病和创伤仅意味着健康的缺失；因为当治愈发生时，并不意味着曾存在的恶——即疾病和创伤——离开身体而移居别处：它们彻底停止存在；因为创伤或疾病并非实体，而是肉体实体的一种缺陷——肉体本身是实体，因此是某种善物，而那些恶——即我们称为健康的善的缺失——是附属于其上的偶性。同样地，灵魂中所谓的恶习，无非是自然之善的缺失。当它们被治愈时，并不转移到他处：当它们在健康的灵魂中停止存在时，它们便不能在任何其他地方存在。\par

% structure: p0030 source=h3 target=subsection reason=relative-heading-rank; base=h2

\subsection{第十二章　一切受造之物原是好的，但因未受造得完美，故易于朽坏。}

因此，所有存在的事物，既然它们的创造者是至善的，它们自身便是善的。但因为它们不像它们的创造者那样是至善且不变地善，它们的善可以减弱和增加。但善的减弱是一种恶，尽管无论减弱多少，若受造物要继续存在，就必须保留某些善以构成该受造物。因为无论受造物何等微小或何种类型，使其成为受造物的善若不摧毁受造物本身便不能毁灭。未朽坏的本性应当受到尊重。但是，若它更进一步是不朽坏的，则无疑被认为具有更高的价值。然而，当它被朽坏时，它的朽坏是一种恶，因为它被剥夺了某种善。因为若不剥夺任何善，它就不受损害；但它确实受到损害，因此它被剥夺了善。所以，只要受造物处于朽坏的过程中，其中就有某种善正被剥夺；如果受造物的一部分应留存而无法被朽坏，那么这必定是一个不朽坏的受造物，因此朽坏的过程将导致这伟大善的彰显。但如果它不停止被朽坏，它也不能停止拥有朽坏可剥夺的善。但如果它被朽坏彻底完全地吞噬，那么就将不再有善留下，因为将不再有受造物。因此，朽坏只能通过吞噬受造物来消耗善。所以，每个受造物都是善的；若是不可朽坏的，便是大善；若是可朽坏的，则是小善；但无论如何，只有愚昧无知的人才会否认它是善的。如果它被朽坏完全吞噬，那么朽坏本身也必停止存在，因为没有受造物留存供它栖身。\par

% structure: p0032 source=h3 target=subsection reason=relative-heading-rank; base=h2

\subsection{第十三章　无善之处即无恶；恶人是恶的善。}

因此，若没有善，就没有我们所谓的恶。但完全无恶的善是完美的善。另一方面，含有恶的善是有缺陷或不完美的善；而没有善之处便不可能有恶。从这一切我们得出一个奇特的结果：既然每个受造物就其为受造物而言是善的，那么当我们说一个有缺陷的受造物是恶的受造物时，我们似乎只是在说善的是恶的，并且除了善之外，没有任何东西可以是恶的，因为每个受造物都是善的，而且恶只能存在于受造物中。所以，除了某种善的事物之外，没有任何事物可以是恶的。虽然这一陈述看似矛盾，但严格的推理使我们无法回避这一结论。然而，我们必须警惕不要招致先知的谴责：\BibleQuote{祸哉，那些称恶为善，称善为恶，以暗为光，以光为暗，以苦为甜，以甜为苦的人！}然而我们的主说：\BibleQuote{恶人从心里所存的恶，就发出恶来。}那么，什么是恶人？不就是恶的受造物吗？因为人是受造物。如果人因是受造物而成为善的事物，那么恶人岂不是恶的善？然而，当我们准确区分这两件事时，我们发现他之所以是恶，并非因为他是人，他之所以是善，也并非因为他邪恶；而是因为他是人所以是善，因为他邪恶所以是恶。因此，谁说\Quote{为人是恶}，或说\Quote{邪恶是善}，就落入了先知的谴责：\BibleQuote{祸哉，那些称恶为善，称善为恶的人！}因为他谴责了天主的创造——即人，却赞扬了人的缺陷——即邪恶。因此，每个受造物，即使是有缺陷的，就其是受造物而言是善的，就其为缺陷而言是恶的。\par

% structure: p0034 source=h3 target=subsection reason=relative-heading-rank; base=h2

\subsection{第十四章　善与恶是相反属性不能同属一主之规则的例外。恶生于善中，且不能离善而存在。}

因此，对于这些我们称为善与恶的相反者而言，逻辑学家的规则——两个相反者不能同时述说同一事物——并不成立。没有天气同时是黑暗和光明的；没有食物或饮料同时是甜和苦的；没有物体在同一时间同一地点是黑和白的；也没有在同一时间同一地点是丑陋和美丽的。而且我们发现这一规则对于许多、实际上几乎所有的相反者都成立，即它们不能同时存在于任何一物中。但是，虽然无人能怀疑善与恶是相反者，它们不仅能够同时存在，而且恶不能脱离善而存在，也不能存在于不是善的事物中。然而，善可以脱离恶而存在。因为一个人或一个天使可以不存在邪恶而存在；但除了人或天使之外，没有事物可以是邪恶的：就他是人或天使而言，他是善的；就他是邪恶而言，他是一种恶。而且这两个相反者如此共存，以至于如果善不存在于恶中，恶也无法存在；因为如果没有任何可以败坏的事物，败坏便既无栖息之地，也无来源可生；没有事物可以被败坏，除非它是善的，因为败坏无非就是善的毁灭。因此，从善中产生了恶，并且除非在善中，恶便不存在；也没有任何其他来源可以让任何恶的性质产生。因为如果有，那么就此物是一个存在而言，它当然是一种善；一个不可败坏的存在会是一种大善；即使是一个可败坏的存在，也必然在某种程度上是善的，因为只有败坏其内的善，败坏才能对它造成伤害。\par

% structure: p0036 source=h3 target=subsection reason=relative-heading-rank; base=h2

\subsection{第十五章 上述论证绝无与主的话牴牾之处：\BibleQuote{好树不能结坏果子}。}

但是，当我们说恶从善而生时，不要让任何人认为这与主的话相矛盾：\BibleQuote{好树不能结坏果子}。因为，正如那真理者所说，你不能从荆棘上摘葡萄，因为葡萄不长在荆棘上。但我们看到，在好土壤上，葡萄树和荆棘都能生长。同样，正如坏树不能结好果子，同样，恶意的意志不能产生善工。但从人的本性（它是善的）中，能产生善的或恶的意志。而且，起初确实没有别的来源让恶的意志产生，除了天使或人的本性（它是善的）之外。而主亲自在他说到树和果子的同一地方明确地表明了这一点。因为他说：\BibleQuote{或使树好，它的果子也好；或使树坏，它的果子也坏}——这足够清楚地警告我们，坏果子不长在好树上，好果子也不长在坏树上；然而，土地本身（他指当时他所教导的那些人）可以生长任何一种树。\par

% structure: p0038 source=h3 target=subsection reason=relative-heading-rank; base=h2

\subsection{第十六章 知道自然界的剧变缘由并非人之幸福所必需；而知道善与恶的缘由则是必需的。}

如今，鉴于这些考虑，当我们喜欢马罗的那句诗：\Quote{幸福的人，他获得了对事物缘由的认识}，我们不应认为知道自然剧变的缘由是幸福所必需的，这些缘由隐藏在自然界王国最隐秘的深处：\Quote{地震的力使深海涌起，冲破其堤坝，然后又回缩、沉降}。但我们应当知道善与恶的缘由，只要人在今生能够知道它们，以便避免今生如此多的错误和烦恼。因为我们的目标始终是达到那样的幸福状态：没有烦恼会折磨我们，没有谬误会误导我们。如果我们必须知道自然界的剧变的缘由，那么与我们健康相关的缘由是我们更当知道的。但鉴于我们对此无知，不得不求助于医生，那么对于我们对于天地间隐藏秘密的无知，我们或许可以相当耐心地承受。\par

% structure: p0040 source=h3 target=subsection reason=relative-heading-rank; base=h2

\subsection{第十七章 错误的性质。并非一切错误都有害，尽管尽可能避免它是人的本分。}

因为尽管我们应当以最大的谨慎避免错误，不仅在大事上，即使在小事上也是如此，而且我们只能由于无知而犯错误，但这并不意味着，如果一个人对某事无知，他就必然会立即陷入错误。那倒是自以为知道其实并不知道之人的命运。因为他把假的当作真的接受，这正是错误的本质。然而，一个人犯错误的主题是什么，这一点非常重要。因为在同一主题上，我们理所当然地更看重一个受过教育的人，而不是一个无知的人；更看重一个没有犯错误的人，而不是一个犯了错误的人。然而，在不同主题的情况下——即一个人知道一件事，另一个人知道另一件事，而且前者的知识是有用的，后者的知识不那么有用，甚至是有害的——谁不会在后者所知的事情上，宁愿选择前者的无知，而不是后者的知识呢？因为有些事情，无知比知识更好。同样，有时候偏离正路也有好处——不过是在旅行中，而不是在道德上。我自己就曾遇到过这种情况，在两条路交汇处走错了路，结果没有经过一伙多纳徒派武装分子埋伏的地方。然而我虽然绕了路，还是到达了我要去的地方；当我听说那埋伏时，我为自己的错误感到庆幸，并为此感谢天主。现在，谁不愿做这样犯了错误的旅行者，而不愿做那个没有犯错误的强盗呢？因此，或许诗人之王让一个痛苦的情人说出这样的话：\Quote{我是多么不幸，我被一个邪恶的错误所掳走！} 因为有一种错误是好的，它不仅毫无害处，而且产生了实际的好处。但当我们更仔细地审视真理的本性，并考虑到犯错误就是把假的当作真的，把真的当作假的，或者把确定的当作不确定的，把不确定的当作确定的，并且当我们在说出来或同意它时，说\Quote{是，是；非，非}，错误在灵魂中越是显得美丽和合理，它就越是丑陋和可憎——那么，我们所过的这种生命确实是悲惨的，仅仅因为这个原因，有时为了保存生命，不得不陷入错误。愿那种生命不是这样，在那里真理本身就是灵魂的生命，在那里没有人欺骗，也没有人被欺骗。但在这里，人欺骗人也被人欺骗，当他们误导他人时，比他们自己因相信说谎者而被误导更值得同情。然而，理性灵魂如此厌恶虚假，如此努力地对抗错误，以至于连那些喜欢欺骗的人也不愿被欺骗。因为说谎者并不认为自己犯错，而是认为自己引导了信任他的他人进入错误。当然，就他所撒谎的内容而言，他并没有错，如果他自己知道真相的话；但他在这件事上受骗了：他认为他的谎言对他无害，然而每一个罪对犯罪者本身的伤害都大于对受害者的伤害。\par

% structure: p0042 source=h3 target=subsection reason=relative-heading-rank; base=h2

\subsection{第十八章 绝不允许说谎；但谎言根据意图和主题在罪责上大不相同。}

但这里出现了一个非常困难且极其复杂的问题，我曾为此写了一本大书，认为必须予以回答。问题是：一个好人是否在任何时候都有义务说谎？因为有些人甚至主张，在有些场合，发虚誓、就敬拜天主的事以及天主本身的属性说假话，也是一件良善虔诚的事。但在我看来，似乎每一谎言都是罪，尽管说谎的意图和主题有很大区别。因为为帮助人而说谎的罪，不如为伤害人而说谎的罪严重；撒谎将旅行者引入歧途的人，其伤害不如以虚假或误导的陈述扭曲人一生历程的人。当然，没有人应被定为说谎者，如果他所说的是假的但他相信是真的，因为这样的人并非故意欺骗，而是自己受了骗。同理，一个人也不应被指控说谎——尽管他有时可能被指责为轻率——如果他因疏忽而接受了虚假并当作真理；反之，一个人所说的是真的但他相信是假的，就其自身意识而言，他是说谎者。因为他说了他不相信的，对他自己的良心来说是假的，尽管事实上可能是真的；一个口里说真话却不知是真是假，但心里愿意说谎的人，绝不脱离说谎。因此，不看所说之事，只看说话者的意图，无意中说谎而始终以为是真的，比无意中说真话而良心却意图欺骗的人要好。因为前者不是想一套说一套；后者虽然陈述事实上可能为真，但心里想的和嘴上说的不一样：这正是说谎的本质。但当我们考虑与所说主题相关的真理与虚假时，一个人欺骗或受骗的点就变得极为重要。因为虽然就人的良心而言，欺骗比受骗是更大的恶，然而，就与宗教无关的事说谎，远比在与正确敬拜天主所必须认识和相信的事上陷入错误，其恶要小得多。举例说明：假设一个人对某个已死的人说他还活着，明知这是假的；另一个人受骗，相信基督会在某个时间（时间长短随你定）死去；那么，像前者那样说谎，岂不比像后者那样受骗好得无可比拟？将某人引入前一种错误，岂不是远比被任何人引入后一种错误危害小得多？\par

% structure: p0044 source=h3 target=subsection reason=relative-heading-rank; base=h2

\subsection{第十九章。人的错误所产生之恶大小极不相同；但每个错误本身都是恶。}

在某些事上，受骗是极大的恶；在某些事上是小恶；在某些事上根本不是恶；在某些事上甚至是实际的益处。当一个人不信导向永生的事，或信导向永死的事时，他受骗是极其有害的。一个人受骗，因以假为真，给自己带来暂时的烦恼，这是小恶；因为信徒的忍耐甚至会将此用于善，例如，当他把坏人当作好人而受到他的伤害时。但一个人相信坏人是好人，却没有受到任何伤害，他并不因受骗而更糟，也不落在先知的谴责之下：\BibleQuote{祸哉，那些称恶为善的人！}因为我们要明白，这话不是论恶人说的，而是论使人为恶的事说的。因此，称奸淫为善的人，正应受先知的谴责。但称奸夫为善，认为他贞洁，而不知道他是奸夫的人，并没有在善恶的性质上犯错误，只是在人的行为的隐秘上犯了错误。他称那个人为善，是凭着他相信那无疑是善的；他称奸夫为恶，称纯洁的人为善；他称这个人为善，是不知道他是奸夫，而相信他是纯洁的。此外，如果因犯错而免于一死，如我以前曾经历过的那样，人甚至从自己的错误中得些益处。但当我断言在某些情况下人可能受骗而毫无损害，甚至得到一些益处时，我并不是说错误本身不是恶，或以任何意义是善；我只是指因犯错而避免的恶或获得的益处。因为错误本身是恶：若关乎大事则是大恶，若关乎小事则是小恶，但总是恶。因为心智健全的人谁能否认，以假为真、以真为假，或以不定为定、以定为不定，是恶呢？但将实际是坏的人当作好人，这是一种错误；而因这一错误未遭受进一步的损害——假设我们认为好的坏人没有对我们造成伤害——则是另一回事。同样，认为自己在正路上而实际不在，这是一种错误；而这一错误（是恶）导致某种好处，如救我们脱离恶人的埋伏，则是另一回事。\par

% structure: p0046 source=h3 target=subsection reason=relative-heading-rank; base=h2

\subsection{第二十章。每个错误并非都是罪。对学院派哲学家意见的检讨，即为避免错误应在一切情形中止判断。}

我不确定以下这类错误是否应当称为罪：例如，对一个坏人形成良好评价而不知其为人；或者，不是通过身体感官的通常知觉，而是出现其他类似表象，我们在精神中感知它们，却以为是在身体中感知，或者在身体中感知却以为是在精神中感知（如同宗徒伯多禄所犯的错误，当天使突然将他从锁链和监禁中释放出来时，他以为自己看见了神视）；或者，就可感觉事物本身而言，我们误将粗糙当作光滑，或将苦当作甜，或以为腐臭之物有好气味；或者，误将马车经过当作雷声；或者误认一人为另一人，两人极为相似，如同双胞胎常见的情况（因此我们伟大的诗人称之为\Quote{使父母喜悦的错误}）——这些错误以及其他类似错误，是否应当称为罪。我现在也不想解决一个非常棘手的问题，这个问题曾困扰那些极其敏锐的思想家，即学院派哲学家：智者是否应当对任何事物表示同意，既然他可能因同意虚假而陷入错误：因为他们断言，一切事物要么未知，要么不确定。在我归化后不久，我写了三卷书，以排除那些仿佛立于信仰门槛上的反对意见。确实，在起初就必须排除这种对达至真理的彻底绝望，这种绝望似乎被这些哲学家的论证所加强。在他们眼中，每个错误都被视为罪，他们认为只有完全暂停信念才能避免错误。因为他们说，同意不确定之事的人会陷入错误；他们用最敏锐但最大胆的论证试图证明，即使一个人的意见偶然为真，其真实性也并非确定，因为无法区分真理与谬误。但在我们这里，\Quote{义人必因信德而生活}。如果同意被取消，信德也随同消失；因为没有同意就不可能有信仰。有些真理，无论我们知道与否，若要获得幸福生活，即永生，就必须相信。但我不知道是否应当与这些人争辩，他们不仅不知道眼前有永生，甚至不知道此刻自己是否活着；不，他们竟说不知道他们不可能不知道的事。因为任何人不可能不知道自己活着，因为如果他不是活的，他就不可能不知道；因为不仅知识，就连无知也只能属于活人。但是，他们却以为通过不承认自己活着就能避免错误，而正是他们的错误本身证明了他们是活的，因为一个不活的人不能犯错。因此，正如我们活着不仅是真实的，而且是确定的，同样还有许多其他事物既是真实又是确定的；天主禁止将拒绝同意这些事物称为智慧，而非极度愚蠢。\par

% structure: p0048 source=h3 target=subsection reason=relative-heading-rank; base=h2

\subsection{第二十一章 错误虽不常是罪，却常是恶。}

至于那些相信与否、乃至其真、假定真或假，对于获得天国毫无重要性的问题，在这类事上犯错不被视为罪，或至少是极小极轻微的罪。简言之，在这类事上犯错，无论其性质与大小，都与接近天主的途径无关，这途径就是那\Quote{藉爱德而工作}的基督信德。因为双胞胎孩子事件中\Quote{使父母喜悦的错误}并未偏离这条途径；宗徒伯多禄也未偏离，当他以为自己看见了神视，以致将一事误认为另一事，直到释放他的天使离开他，他仍未能区分他在其中活动的真实对象与梦中的幻象；先祖雅各伯也未偏离，当他相信他实际上活着的儿子已被野兽吞噬。在这些以及其他同类错误中，我们确实被欺骗，但我们对于天主的信德仍稳固。我们迷失，但我们并未离开通往祂的道路。然而，这些错误虽不是罪，却应算作今生的种种恶之一，今生如此受制于虚幻，以至我们将虚假当作真实接纳，拒绝真实如同虚假，紧抓不确定如同确定。尽管它们并不触及那使我们达到永福的真实而确定的信德，却与我们目前所处的悲惨境遇密切相关。确实，若我们如今就享有那未来的真实而完美的幸福，我们就决不会通过任何身体或精神的感官而受任何欺骗。\par

% structure: p0050 source=h3 target=subsection reason=relative-heading-rank; base=h2

\subsection{第二十二章 即使为救他人免受伤害，也不可说谎。}

然而，一切谎言都当称为罪，因为人不仅在他认识真理时，甚至在他可能犯错或受欺骗时，都有责任说出心中所想，无论那想法是真实，还是他只以为是真实。但每一个说谎者都存心欺骗，说出与心中所想相反的话。显然，言语赐予人，不是要人借此彼此欺骗，而是使人能将自己的思想传达给他人。因此，为欺骗的目的而使用言语，而不是为了它被赋予的宗旨，便是罪。我们也不可认为，有时借着说谎能有利于他人，就说有某种谎言不是罪。因为也能借着偷窃来做同样的事，比如我们从富户那里偷取他永不觉察的损失，施舍给穷人，使他因所得而明显受益。同样也可以说奸淫，比如某位妇人似乎会因爱情而死，除非我们满足她的愿望；而若她活着，她本可藉悔改洁净自己。然而无人会因此断言这样的奸淫不是罪。如果我们对贞洁如此重视，那么我们对真理又有什么冒犯，以至于尽管无望为他人谋利，我们也不愿因奸淫而违背贞洁，却愿意因说谎而违背真理？不可否认，那些除非为救人免于伤害否则永不说谎的人，已达至很高的良善标准；但在这到这一标准的人身上，应受称赞的并非欺骗，而是其善意，有时甚至应受奖赏。欺骗应被宽容已足够了，而不应成为赞美对象，尤其是在新盟约的承继者中间，经上对他们说：\BibleQuote{你们的话，是就说是，不是就说不是；除此之外，都是出于邪恶。}\BibleRef{玛 5:37}正因这邪恶——当我们还穿着这必死的肉身时，它从不停止侵入——基督的同承继者自己也说：\BibleQuote{宽免我们的罪债。}\BibleRef{玛 6:12}\par

% structure: p0052 source=h3 target=subsection reason=relative-heading-rank; base=h2

\subsection{第二十三章　先前讨论结果的概要}

既然我们应当认识善恶的原因，至少认识那些足以引导我们走向天国的路上所需的部分，在那里将有不受死亡阴影的生命、没有错误掺杂的真理、没有忧伤侵扰的幸福，我已以我有限篇幅所要求的简洁讨论了这些主题。我认为如今已无可怀疑：我们所享的一切善的唯一原因是天主的良善，而恶的唯一原因是受造为良善却可变者——首先是一位天使，然后是人的坠落，离开那不变的善。\par

% structure: p0054 source=h3 target=subsection reason=relative-heading-rank; base=h2

\subsection{第二十四章　恶的次要原因是无知和私欲}

这就是降临于理性受造物身上的第一个恶——即它第一次丧失善。随之潜入的是，甚至现在违反人的意愿，\newline{}\Emph{无知}于本分，以及\Emph{私欲}追求有害之物；这些又带来了\Emph{错误}和\Emph{痛苦}，当这些感觉临近时，便产生那称为\Emph{恐惧}的心灵退缩。此外，当心灵达到其欲望的目标时，无论这些目标何等有害或空虚，错误阻止它认识其真实性质，或由于病态的欲望压倒了它的认知，以致它因\Emph{愚昧的快乐}而自傲。从这些恶的泉源——它们源于缺陷而非过剩——流出困扰理性本性的每一种苦难形式。\par

% structure: p0056 source=h3 target=subsection reason=relative-heading-rank; base=h2

\subsection{第二十五章　天主对堕落的人和天使的审判。身体的死亡是人所特有的惩罚}

然而，这样的本性，在一切恶中，仍不能失去对幸福的渴望。现在我提到的这些恶，是共有的，适用于所有因自己的邪恶而被天主公正定罪者，无论他们是人还是天使。但有一种惩罚是人所特有的——身体的死亡。天主曾以死亡的惩罚警告他，若他犯罪，确实让他自由运用自己的意志，但命令他服从，违者受死；祂将他置于伊甸的幸福中，如同在生命的庇护所中，意图是：若他保持义，就从那里上升到一个更好的地方。\par

% structure: p0058 source=h3 target=subsection reason=relative-heading-rank; base=h2

\subsection{第二十六章　因亚当的罪，他整个后裔都败坏了，并生来即承受他所招致的死亡惩罚}

此后，他犯罪后被逐流放，因他的罪，整个以他为根的后裔在他里面败坏了，并因而受死亡惩罚的辖制。因此，所有从他与那引诱他犯罪并与他一同被定罪的女人所生的后代——这后代是情欲的产物，同一种不服从的惩罚也加于其上——都沾染了原罪，并借此经由各种错误和痛苦，被拖入那最后的无尽惩罚，他们与堕落的天使（他们的败坏者、主人和同样命运的分享者）共同承受这惩罚。因此，\newline{}\BibleQuote{藉着一个人，罪恶进入了世界，藉着罪恶，死亡进入了世界；这样，死亡就传递到所有人身上，因为所有人都犯了罪。}\BibleRef{罗 5:12}宗徒在此处所说的\Quote{世界}，当然是指全人类。\par

% structure: p0060 source=h3 target=subsection reason=relative-heading-rank; base=h2

\subsection{第二十七章　亚当之罪使人类陷入的苦难状态，以及藉天主的慈悲实现的恢复}

如此，情势便已定。整个人类都处在定谳之下，深陷于苦难中，从一种邪恶被抛向另一种，且与堕落的天使同谋，为那邪恶的反叛而承受应得的惩罚。因为恶者自由地以盲目且无度的私欲所行的一切，以及他们违愿地作为公开惩罚所受的一切，显然都属于天主的义怒。然而，造物主的美善从未间断：祂供给恶天使生命与活力（没有这些，他们的存在将即刻终结）；对于出自受咒诅败坏的种族的人类，祂赋予其种子以形态与生命，塑造他们的肢体，并透过他们生命的各个时期、在地球的不同角落，激发他们的感官，赐予他们所需的滋养。因为祂判定，从恶中引出善，比不容许任何恶存在更好。若祂决定，对于人类，如同对堕落的天使一样，不应有幸福的重建，那么，那反叛天主、滥用自由轻蔑并违犯造物主诫命（本可轻易遵守）、固执地背离光明而玷污自身内造物主的肖像、恶用自由意志脱离造物主律法之健全束缚的受造物——难道完全永远被天主遗弃、并承受其应得的永罚，不正是公义的吗？的确，若天主仅是公义而非仁慈，若祂未定意使祂不应得的仁慈在与受惠者之不堪的对比中更显光辉，祂便会如此行。\par

% structure: p0062 source=h3 target=subsection reason=relative-heading-rank; base=h2

\subsection{第28章. 当反叛的天使被逐出时，其余的天使与天主同在永恒幸福中安然享受。}

因此，当一些天使在其骄傲与不虔中反叛天主，并从其天堂居所被投下至最深黑暗时，其余的天使则与天主同在永恒不变之纯洁与幸福中安居。因为他们并非全部出自一位堕落受罚的天使，以致像人类一样，由一项原罪牵连于遗传的罪债，因而承受那一位所招致的惩罚；而是，当那后来成为魔鬼者及其同谋之罪在骄傲中被高举，且因那高举而与他们一同被投下时，其余的天使仍坚定地忠于对主的虔诚与服从，并获得了他们先前未曾享有的、关于其永恒安全及免于堕落之可能性的确知。\par

% structure: p0064 source=h3 target=subsection reason=relative-heading-rank; base=h2

\subsection{第29章. 人类被恢复的部分将依天主的许诺，继承反叛天使所失的地位。}

因此，宇宙的造物主及治理者天主乐意如此：既然全体天使并未都陷入反叛，那一部分堕落的天使应永受丧亡，而另一部分在反叛中坚定忠诚的天使，应为其永恒幸福的确知而喜乐；另一方面，人类，作为理智受造物中的剩余部分，既全无例外地因原罪及本罪以及随之而来的惩罚而丧亡，应有一部分被恢复，并填补魔鬼反叛与堕落所留下的天使团体中的空缺。因为给圣徒的应许是：在复活时，他们将等同于天使。如此，那在上面的耶路撒冷，即我们众人的母亲、天主的城，将不会丧失其任何一位公民，或许还将统治更为丰裕的人口。我们不知道圣徒或魔鬼的数目；但我们知道，那被称为在地上不育的圣母亲的孩子，将继承堕落天使的位置，并永远居住于他们堕落后的那个平安之所。但公民的数目，无论是现今的或是将来的，都呈现在伟大造物主的思想中，祂呼唤那不存在的如同存在一般，并以度量、数目和重量安排万物。\par

% structure: p0066 source=h3 target=subsection reason=relative-heading-rank; base=h2

\subsection{第30章. 人得救并非因善行，也非因自己意志的自由决定，而是因着天主的恩宠，藉着信德。}

然而，天主应许赦免并分享祂永恒国度的这部分人类，他们能凭借自己行为的功绩而恢复吗？断乎不可。因为一个失丧的人能行什么善行，除非他已经从灭亡中被释放？他们能凭自己自由意志的决定行任何事吗？我再说，断乎不可。因为人因滥用其自由意志而既毁灭了它，也毁灭了自己。因为正如一个自杀的人，他在自杀时当然活着，但自杀后就不再存活，也不能使自己复生；同样，当人凭自己的自由意志犯罪，罪恶战胜了他，他的意志自由就丧失了。\Quote{因为人被谁制伏，就是谁的奴仆。}这是宗徒\Person{伯多禄}的判断。既然这确实是真的，我问，这奴仆能拥有怎样的自由，除非在他乐意犯罪的时候？因为甘心乐意行主人旨意的人，是自由地受奴役。因此，犯罪的人有犯罪的自由。所以，他不会有行善的自由，直到他从罪恶中得到释放，开始成为正义的奴仆。这才是真正的自由，因为他对正义的作为感到欣喜；同时这也是神圣的奴役，因为他服从天主的旨意。然而，这行善的自由从何而来，给那被卖在罪恶之下的人？除非他被那一位救赎，祂曾说过：\Quote{如果圣子使你们自由，你们就必得以自由。}在人在内完成这种救赎之前，当他还没有行善的自由时，他怎能谈论他的自由意志和善行呢？除非他被那愚昧的骄傲夸口所膨胀，而宗徒制止了这种骄傲，他说：\Quote{你们得救是因着恩宠，借着信德。}\par

% structure: p0068 source=h3 target=subsection reason=relative-heading-rank; base=h2

\subsection{三十一、信德本身是天主的恩赐；凡相信的人必不缺善行。}

并且，免得人们至少将自已的信德的功绩归给自己，不明白这同样是天主的恩赐，这位同一位宗徒，在另一处说他\Quote{蒙了天主的仁慈，得以成为忠信的}，在此也加上：\Quote{并不是出于你们自己，而是天主的恩赐：不是出于行为，免得有人自夸。}并且，免得有人认为相信的人会缺乏善行，他进一步说：\Quote{因为我们原是祂的化工，在基督耶稣内造成的，为行善，就是天主预先预备好要我们行在其中的。}这样，我们才能真正获得自由，当天主塑造我们，即重新形成并创造我们，不是作为人——因为祂已经做了那事——而是作为善人，这是祂的恩宠现在正在做的，使我们成为在基督耶稣内的新受造物，正如经上所说：\Quote{天主，求祢为我再造一颗洁净的心。}因为天主已经创造了我的心，就人类心脏的物理结构而言；但圣咏作者祈求的是那仍在他心中存留的生命得以更新。\par

% structure: p0070 source=h3 target=subsection reason=relative-heading-rank; base=h2

\subsection{三十二、意志的自由也是天主的恩赐，因为天主在我们内运作，使我们愿意并实行。}

并且，若有人倾向于夸口，不是夸自己的行为，而是夸自己意志的自由，仿佛最初的功绩属于他，这行善的自由本身是作为他应得的赏报赐给他的，让他听这同一位恩宠的宣讲者的话，他说：\Quote{因为是天主在你们内运作，使你们愿意并实行，为成就祂的美意；}在另一处又说：\Quote{这样看来，它并不在于愿意的人，也不在于奔跑的人，而在于显示仁慈的天主。}然而，无疑地，如果一个人到了能运用理性的年龄，他除非愿意，就不能信、望、爱，也不能获得天主至高召叫的奖赏，除非他自愿奔跑；那么在什么意义上说\Quote{并不在于愿意的人，也不在于奔跑的人，而在于显示仁慈的天主}，除非正如经上所记：\Quote{心的准备在于主？}否则，若说\Quote{并不在于愿意的人，也不在于奔跑的人，而在于显示仁慈的天主}，是因为它在于两者，即既在于人的意愿也在于天主的仁慈，以至我们应当理解这话\Quote{并不在于愿意的人，也不在于奔跑的人，而在于显示仁慈的天主}似乎意味着单凭人的意愿是不够的，除非天主的仁慈随之——那么就会推论出，单凭天主的仁慈是不够的，除非人的意愿随之；因此，我们若可以正确地说\Quote{并不在于愿意的人，而在于显示仁慈的天主}，因为人的意愿本身不够，为什么我们不可以同样正确地反过来说：\Quote{并不在于显示仁慈的天主，而在于愿意的人}，因为天主的仁慈本身并不足够？当然，如果没有任何基督徒敢说\Quote{并不在于显示仁慈的天主，而在于愿意的人}，免得他公开反对宗徒，那么真正的解释是：\Quote{并不在于愿意的人，也不在于奔跑的人，而在于显示仁慈的天主}全在于天主，祂既使人的意志成为正义，从而预备它接受援助，又援助它当它预备好了。因为人的意志的正义先于天主的许多恩赐，但并非所有；并且它本身必须被包括在它所不先于的那些恩赐之中。我们在圣经中读到，天主的仁慈\Quote{要迎接我}，祂的仁慈\Quote{要追随我}。它走在不愿者之前，使他愿意；它跟随愿意者，使他的意愿有效。为什么我们被教导要为仇敌祈祷，他们显然不愿意过圣善的生活？除非天主在他们内运作意愿。为什么我们自己也受教导祈求以便领受？除非那在我们内创造了愿望的那一位，祂亲自满足那愿望。因此，我们为仇敌祈祷，使天主的仁慈在他们之先，如同它在我们之先一样；我们为自己祈祷，使祂的仁慈追随我们。\par

% structure: p0072 source=h3 target=subsection reason=relative-heading-rank; base=h2

\subsection{三十三、人，生来是忿怒之子，需要中保。天主在何种意义上被称为发怒。}

然而，人类处于公义的定罪之下，所有人都成了可怒之子。关于这忿怒，经上记载说：\BibleQuote{我们度尽的年岁都在你的忿怒中；我们度过的年岁好像一声叹息。} 关于这忿怒，约伯也说：\BibleQuote{妇人所生的人，日子短少，满有愁烦。} 关于这忿怒，主耶稣也说：\BibleQuote{信子的人有永生；不信子的人不得见永生，天主的震怒却常在他身上。} 祂不是说震怒将要临到，而是说\BibleQuote{常在他身上。} 因为每个人生来就带着这忿怒；因此宗徒说：\BibleQuote{我们本为可怒之子，和别人一样。} 当人类因原罪处于这忿怒之下，并且这原罪因所加的本罪的数量和严重程度而更加沉重和致命时，就需要一位中保，即一位和好者，祂藉着奉献一次祭献（法律和先知的一切祭献都是这祭献的预像）来除去这忿怒。因此宗徒说：\BibleQuote{因为我们作仇敌的时候，且藉着天主圣子的死，得与天主和好；既已和好，就更要因祂的生命得救了。} 当说天主发怒时，我们并非认为祂有人类愤怒时的那种激动情绪；而是藉着类比从人的情感借用\Quote{愤怒}一词来称呼祂对罪的公义不悦。但我们藉着中保与天主和好，并领受圣神，使我们这些仇敌成为儿女（\BibleQuote{凡被圣神引导的，都是天主的子女}），这就是藉着我们的主耶稣基督而来的天主的恩宠。\par

% structure: p0074 source=h3 target=subsection reason=relative-heading-rank; base=h2

\subsection{第34章 中保基督藉童贞玛利亚降生成人的不可言喻的奥秘。}

如今，关于这位中保，若要说得与祂相称，篇幅太长；而且，说与祂相称的话，并非人力所能及。因为谁能用连贯的话语解释这一句话：\BibleQuote{圣言成了血肉，居住在我们中间}，使我们相信全能天主圣父的独生子，因圣神和童贞玛利亚而生。圣言成为血肉的意思，并不是天主性改变了血肉，而是天主性取了我们的血肉。我们在这里理解\Quote{血肉}为\Quote{人}，以部分代全体，正如经上说：\BibleQuote{凡有血气的，没有一个因行律法称义}，即没有人。因为我们应当相信，祂所取的人性无任何缺乏，只是这人性完全不受任何罪染——不是像那通过两性肉体情欲所孕育的、出生在罪中、其罪债在再生中被洗净的本性；而是适合由童贞女所生，那时母亲的信仰，而非情欲，是怀孕的条件。如果她的童贞在分娩时受到损害，祂就不是由童贞女所生了；那么祂由童贞玛利亚所生就是假的（愿天主禁止），而这是整个教会所信和宣认的，教会效仿祂的母亲，每日生产祂身体的肢体，却仍然保持童贞。若你愿意，请阅读我写给那位杰出人士——我怀着敬意和感情提及他的名——\Person{沃卢西亚努斯}的关于圣玛利亚童贞的信函。\par

% structure: p0076 source=h3 target=subsection reason=relative-heading-rank; base=h2

\subsection{第35章 耶稣基督，既是天主的独生子，同时又是人。}

因此，基督耶稣、天主子，既是天主又是人；在万世之前是天主，在我们的世界是人：祂是天主，因为祂是圣言（因为\BibleQuote{圣言就是天主}）；祂是人，因为在祂同一身位中，圣言与一个身体和一个理性灵魂结合。因此，就祂是天主而言，祂与圣父原是一体；就祂是人而言，圣父比祂大。因为祂本是天主的独生子，不是因恩宠，而是因本性，为要也成为充满恩宠的，祂成了人子；祂在自己的同一性中结合了两性，两性构成一位基督；因为\BibleQuote{祂虽具有天主的形体，却不以自己与天主同等为强夺的}，祂本来与天主同等，\BibleQuote{反而空虚自己，取了奴仆的形体}。祂并没有失去或减少天主的形体。因此，祂同时被贬抑和保持同等，两者合而为一，如已所述：但祂作为圣言是一方面，作为人是另一方面。作为圣言，祂与圣父同等；作为人，祂小于圣父。一位天主子，同时是人子；一位人子，同时是天主子；并非两位天主子，一位是天主一位是人，而是一位天主子：无始的天主，有始的人，我们的主耶稣基督。\par

% structure: p0078 source=h3 target=subsection reason=relative-heading-rank; base=h2

\subsection{第36章 天主的恩宠在提升人基督耶稣至天主子尊严的事上清楚地、显著地显明出来。}

现在，天主的恩宠在此以最大的能力和清晰度显示出来。因为基督内的人性曾有过什么功德，能使它如此空前地被提升到天主独生子位格的合一之中？曾有什么意志的良善、愿望与意向的良善、什么善工先行，使这个人配成为与天主同一位格？难道祂先前曾是个人，然后赢得了这史无前例的赏报，即被认为配成为天主？当然不是；从祂开始成为人的那一刻起，祂就无非是天主子、天主的独生子、成为血肉的圣言，因此祂是天主，以至于正如每个人将身体和理性灵魂合一于一个位格，同样基督将圣言和人合一于一个位格。那么，为什么这种闻所未闻的荣耀被赐予人性——这荣耀既然没有前在的功德，当然完全是恩宠——除非是为了让那些冷静而诚实地看待此事的人，能清楚看见天主白白恩宠的能力，并明白他们正是藉着使基督耶稣之人免于犯罪可能的同一恩宠，而从自己的罪恶中被成义？因此，天使在向基督的母亲预报祂的诞生时，这样问候她：\BibleQuote{万福，满被恩宠的！}稍后又说：\BibleQuote{你在天主前获得了恩宠。}她之所以被称为满被恩宠并在天主前获得了恩宠，是因为她将成为她主的母亲，不，将成为所有血肉之主的母亲。但是，关于基督自己，福音作者若望在说了\BibleQuote{圣言成了血肉，居住在我们中间}之后，又加上：\BibleQuote{我们见了祂的光荣，如同父独生者的光荣，满溢恩宠和真理。}当他说\BibleQuote{圣言成了血肉}时，这就是\BibleQuote{满被恩宠}；当他说\BibleQuote{父独生者的光荣}时，这就是\BibleQuote{满溢真理}。因为真理本身，祂是父的独生子，不是由于恩宠，而是由于本性，祂因恩宠接受了我们的人性，并如此与自己位格结合，以至于祂自己也成了人子。\par

% structure: p0080 source=h3 target=subsection reason=relative-heading-rank; base=h2

\subsection{第三十七章 同一恩宠更清楚显示于此：基督按肉身是由圣神所生。}

因为同一耶稣基督，祂是独生子，即天主的独生子、我们的主，是由圣神和童贞玛利亚所生。我们知道，圣神是天主的恩赐，这恩赐本身的确与施予者平等。因此圣神也是天主，不逊于圣父和圣子。所以，基督在其人性中的诞生是由圣神而生的这一事实，是恩宠的另一个清楚显示。因为当童贞女问天使她所预报的如何成就，因为她不认男人时，天使回答说：\BibleQuote{圣神要临于你，至高者的能力要庇荫你，因此那要诞生的圣者，将称为天主子。}当若瑟想要暗暗休退她，怀疑她犯奸淫，因为他知道她不是由自己怀孕时，天使告诉他：\BibleQuote{不要怕娶你的妻子玛利亚，因为那在她内受孕的是出于圣神。}也就是说，你所怀疑的由别人所生的，是出于圣神。\par

% structure: p0082 source=h3 target=subsection reason=relative-heading-rank; base=h2

\subsection{第三十八章 耶稣基督按肉身不是由圣神所生，以致圣神是祂的父亲。}

然而，我们是否因此要说圣神是基督之人的父亲，并且如同天主圣父生了圣言，同样天主圣神生了这个人，而这两种本性构成一位基督；并且作为圣言，祂是天主圣父之子，而作为人，祂是天主圣神之子，因为圣神作为祂的父亲由童贞玛利亚生了祂？谁敢这样说？也无需用推理来显示这个假设会产生多少其他荒谬之处，因为它本身就是如此荒谬，以致没有信友的耳朵能忍受听到它。因此，我们宣认：\Quote{我们的主耶稣基督，祂出自天主是天主，作为人是由圣神和童贞玛利亚所生，兼有神性和人性两性，是全能天主圣父的独生子，圣神由祂所发。}那么，如果圣神没有生祂，我们在什么意义上说基督是由圣神所生？是不是说圣神造了祂？因为我们的主耶稣基督，虽然作为天主\BibleQuote{万物是藉着祂造成的}，但作为人，祂自己却是受造的；正如宗徒所说：\BibleQuote{祂按肉身是由达味的后裔所生。}但既然童贞女所怀孕和生产的那个受造物（虽然只与圣子的位格结合）是由整个圣三所造（因为圣三的工程是不可分的），为什么唯独圣神被说成是造了它？是不是因为当三位中的一位被提及为任何工程的作者时，应理解为整个圣三在工作？那是真的，并可以用例子证明。但我们无需在此停留更久。因为难题是，在什么意义上说\Quote{由圣神所生}，而祂绝不是圣神之子？因为虽然天主造了这个世界，但说它是天主子或由天主所生就不对了；我们会说它是被天主所创造、所造、所形成、所安排，或任何我们可以恰当使用的表达。那么，当我们宣认基督是由圣神和童贞玛利亚所生时，很难解释为什么祂不是圣神之子，却是童贞玛利亚之子，既然祂是由祂和她所生。毫无疑问，祂由圣神所生，不是如同由父亲所生，与祂由童贞女所生如同由母亲所生，不是同一个意义。\par

% structure: p0084 source=h3 target=subsection reason=relative-heading-rank; base=h2

\subsection{第三十九章 并非一切由另一者所生的，都应被称为那另一者的儿子。}

我们不必因此假定，凡从某物所生的，即刻就可称为那物的儿子。因为，且不谈人子出生与一根头发或虱子从人生出意义不同——后者皆非儿子——姑且略过这点，作为如此重要主题的过于卑微的比喻：可以肯定，那些由水和圣神所生的人，不可恰当地称为水的儿子，虽然他们被称为天主圣父和教会母亲的儿子。同样，那由圣神所生的，乃是天主圣父的儿子，而非圣神的儿子。因为我对头发及其他事物的论述足以向我们显示：并非凡从另一物所生的，皆可称为那所从出之物的儿子；同样，并非所有被称为某人之子的人，都是他所生，因为有些儿子是被收养的。有些人被称为地狱之子，并非因生于地狱，而是因预备为之，如同天国之子预备进入天国一样。\par

% structure: p0086 source=h3 target=subsection reason=relative-heading-rank; base=h2

\subsection{第四十章 基督藉圣神降生，向我们彰显天主的恩宠}

因此，既然一物可由另一物所生，却未必是其儿子；且并非所有称为儿子的，都是被称之父所生，那么，这一安排——基督由圣神所生，却非圣神之子，而由童贞玛利亚所生，作为她的儿子——显然是为显示天主的恩宠。因为这恩宠使一个人，毫无任何先前的功绩，在其作为人之存在的开始，就与天主的圣言如此结合于一位格之中，以致那既是人子者，同时又是天主之子；那既是天主之子者，同时又是人子。在祂的人性被接纳进入神性时，这恩宠本身对那人而言变得如此自然，以致没有给罪恶的进入留下余地。因此，这恩宠由圣神所象征；因为圣神，虽其本性是天主，也可称为天主的恩赐。要足够解释这一切——如果真能完成的话——需要非常冗长的讨论。\par

% structure: p0088 source=h3 target=subsection reason=relative-heading-rank; base=h2

\subsection{第四十一章 基督自身无罪，却为我们成了罪，好使我们与天主和好}

祂既被孕而生，没有肉欲的任何放纵，因此不带原罪；又由天主的恩宠，以奇妙而不可言喻的方式，与圣父的独生子、那按本性而非按恩宠为子的圣言结合于一位格之中，所以自身无罪；然而，由于祂是以罪恶肉身的形状而来，故被称为罪，好被祭献以洗除罪恶。因为在旧盟约之下，为罪所献的祭品被称为罪。而凡这些祭品所预表和影射的那一位，祂本身实在成了罪。因此，宗徒说了\Quote{我们替基督求你们：与天主和好吧}之后，立即加上：\Quote{因为祂使那不知罪的，替我们成了罪，好叫我们在祂内成为天主的正义}。祂没有说，如同一些不正确的抄本所读的，\Quote{那不知罪的，为我们犯了罪}，仿佛基督本人为我们犯了罪；而是说：\Quote{祂使那不知罪的}，即基督、天主，我们当与之和好的那一位，\Quote{替我们成了罪}。因此，祂成了罪，正如我们成了正义（我们的正义不是自己的，而是天主的，不在我们自身，而在祂内）；祂成了罪——不是祂自己的，而是我们的，不在祂自身，而在我们内——藉着祂被钉于其上的罪恶肉身的形状，表明虽然罪不在祂内，但在某种意义上祂向罪死了，因祂死于那罪恶之像的肉身；并且，虽然祂自己从未活过罪恶的旧生命，但藉着祂的复活，祂预表了那从罪恶的旧死中萌生的新生。\par

% structure: p0090 source=h3 target=subsection reason=relative-heading-rank; base=h2

\subsection{第四十二章 圣洗圣事象征我们与基督同死于罪，并与祂同复活于新生}

这就是我们之中所举行的伟大圣洗圣事的含义：凡达到这恩宠的人，都应死于罪，正如祂被称为死于罪，因为祂死于那罪恶之像的肉身；并从重生之泉中升起，如同祂从坟墓中复活，在圣神内开始新生命，无论肉身年龄几何？\par

% structure: p0092 source=h3 target=subsection reason=relative-heading-rank; base=h2

\subsection{第四十三章 圣洗及其所象征的恩宠为所有人敞开，包括婴儿和成人}

因为从初生婴儿到年迈老人，既然没有一人被摒于圣洗之外，同样，没有一人在圣洗中不向罪而死。但婴儿只死于原罪；年长者则也死于其邪恶生活所加于与生俱来之罪上的一切罪恶。\par

% structure: p0094 source=h3 target=subsection reason=relative-heading-rank; base=h2

\subsection{第四十四章 论及罪时，单数常代复数，复数常代单数}

但是，就连后者也常被称为死于罪，虽然他们无疑不是死于某一个罪，而是死于他们在思想、言语或行为上所犯的一切众多实际之罪：因为单数常被用来代替复数，就如诗人所言：\Quote{他们以武装士兵填满其腹}，虽然此处所指的案例中涉及许多士兵。我们在自己的圣经中也读到：\BibleQuote{求主除掉我们身上的蛇。} 他没有说\Emph{众蛇}，尽管民众正遭受许多蛇的侵害；其他情况也是如此。另一方面，当原罪以复数表达时，例如我们说婴儿受圣洗是为赦免\Emph{众罪}，而不说为赦免\Emph{罪}，这是相反的修辞手法，用复数代替单数；如在福音中论及黑落德的死时说：\BibleQuote{因为那些谋害孩子性命的人已经死了}，而不说\Quote{他死了}。在出谷纪中，梅瑟说：\BibleQuote{他们为自己造了金神像}，虽然他们只造了一只牛犊，并指着它说：\BibleQuote{以色列，这就是你的神，领你出埃及地的神}——这里也是用复数代替单数。\par

% structure: p0096 source=h3 target=subsection reason=relative-heading-rank; base=h2

\subsection{第四十五章  在亚当的第一个罪中，包含了许多种罪。}

然而，即使是那一个罪，\BibleQuote{藉着一个人进入了世界，并因此传给了众人}，且因这罪，婴儿也受圣洗，若将其分解为各个元素，就能看出其中包含许多不同的罪。其中有骄傲，因为人选择受自己管辖，而不受天主管辖；有亵渎，因为他不相信天主；有杀人，因为他给自己带来了死亡；有属灵的邪淫，因为人类灵魂的纯洁被蛇的诱惑所腐化；有偷窃，因为他擅自取用了被禁止触摸的食物；有贪婪，因为他贪求超出他应得之物；以及其他任何在仔细反省中能发现包含在这一公认的罪中的罪。\par

% structure: p0098 source=h3 target=subsection reason=relative-heading-rank; base=h2

\subsection{第四十六章  婴儿不仅可能承担第一对父母的罪债，也可能承担其亲生父母的罪债，这是很可能的。}

而且，有一种很有可能的说法是：婴儿不仅承担第一对父母的罪债，也承担其亲生父母的罪债。因为那神圣的审判：\BibleQuote{我必追讨父亲的罪债，归于儿女}，在他们藉着重生进入新盟约之前，确实适用于他们。而厄则克耳所预言的新盟约正是：\BibleQuote{儿子必不承担父亲的罪债，在以色列中也不再流传这俗语：\InnerQuote{父亲吃了酸葡萄，儿子的牙酸倒了。}} 因此，人人必须重生，才能从他生来就有的罪中获得自由。因为后来所犯的罪可以通过补赎来医治，正如我们在圣洗后所见。所以，除非第一次出生是有罪的，且如此有罪，以至于即使是合法婚生的人也说：\BibleQuote{我是在罪孽中成形的，在罪中我母亲怀了我}，否则就不会设立重生。他没有说\Emph{罪孽}或\Emph{罪}，虽然他可以正确地这样说；但他更喜欢说\Quote{罪孽}和\Quote{罪}，因为在那一个传给众人的罪中，正如我上文所示，可以区分出许多罪；此外，还有亲生父母的其他罪，这些罪虽然对改变本性没有同样效力，但除非天主的恩宠和慈悲干预拯救，否则会使儿女承担罪债。\par

% structure: p0100 source=h3 target=subsection reason=relative-heading-rank; base=h2

\subsection{第四十七章  一个人的其他祖先的罪是否归算于他，这是难以决定的。}

但是，关于从亚当到一个人的亲生父母之间其他祖先的罪，很可能会有疑问：每个所生的人是否都承担他们所有累积的恶行，所有倍增的原罪债，以致出生越晚，情况越糟？或者，天主威胁要追讨父亲的罪债直到三、四代，是否因为祂的慈悲没有将祂对祖先之罪的忿怒延伸到更远，以免那些没有得到重生恩宠的人，如果被迫承担从人类起初以来所有祖先的一切罪债并支付应得的惩罚，会被过于沉重的负担压垮？或者，是否能通过更勤奋的查考和更仔细的解释在圣经中找到这个重大问题的其他解决方案？我不敢冒然断言。\par

% structure: p0102 source=h3 target=subsection reason=relative-heading-rank; base=h2

\subsection{第四十八章  第一个罪的罪债如此之大，只能在中保耶稣基督的血中洗去。}

然而，那个罪，在如此完美的幸福统治的地方犯下，具有如此严重的性质，以致整个人类在一个人身上原初地、可以说根本地被定罪了；它只能通过天主与人之间唯一的中保，那人基督耶稣来赦免和涂抹，祂唯独有力量如此出生，以至于不需要第二次出生。\par

% structure: p0104 source=h3 target=subsection reason=relative-heading-rank; base=h2

\subsection{第四十九章  基督在若翰的圣洗中没有重生，而是接受圣洗以给我们谦卑的榜样，正如祂接受死亡，不是作为罪的惩罚，而是为除去世界的罪。}

当时，那些受\Person{若翰}圣洗的人——\Person{基督}本人也受了此洗——并没有重生；但他们通过祂的前驱所从事的职务得到了预备，这前驱呼喊说：\Quote{预备主的道}，为那唯一能令他们重生的一位。因为祂的圣洗不只用水，如\Person{若翰}的圣洗，而是也用\OriginalTerm{圣神}{Holy Ghost}；因此，凡信\Person{基督}的，都由那\OriginalTerm{圣神}{Holy Ghost}重生；\Person{基督}由那\OriginalTerm{圣神}{Holy Ghost}所生，祂不需要重生。因此，祂受洗后所听到的圣父的宣告：\Quote{今日我生了祢}，并非指祂受洗的那一天，而是指永恒不变的一日，显示这个人（耶稣）与\OriginalTerm{独生子}{Only-begotten}是同一性体。因为一个既不以昨日结束为开始，也不以明日开始为结束的日子，是永恒的今日。所以，祂请求\Person{若翰}用水为祂施圣洗，并非要洗去祂的任何不义，而是为彰显祂谦卑的深度。因为圣洗在祂身上找不到任何可洗去的，正如死亡在祂身上找不到任何可惩罚的；这样，魔鬼被粉碎并征服，是基于最严格的正义，而非仅仅凭借暴力；因为，魔鬼极不公正地将\Person{基督}置于死地，尽管\Person{基督}身上没有该死的罪，所以最公义的是\Person{基督}应使魔鬼失去对那些因罪而正义地受制于其捆绑之人的控制。因此，这两者，即圣洗和死亡，都是祂所承受的，并非出于可怜的必然性，而是出于祂对我们自由自愿的怜悯，并作为一项安排的一部分，正如一人将罪带入世界，即带入全人类，同样，一人要除去世界的罪。\par

% structure: p0106 source=h3 target=subsection reason=relative-heading-rank; base=h2

\subsection{第50章 基督不仅除去了那一个原罪，也除去了所有附加的罪。}

其区别在于：第一个人将一个罪带入世界，但此人不仅除去了那一个罪，还除去了祂发现附加在其上的所有罪。因此，宗徒说：\BibleQuote{但恩赐不同于过犯；因为审判是由一人而定罪，恩赐却是由许多过犯而成义。}因为很明显，我们凭本性带来的那一个罪，即使孤立存在，也会使我们被定罪；但恩赐使人从许多过犯中成义：因为每个人，除了与所有同类共同凭本性带来的那一个罪以外，还犯了许多严格属于他自己的罪。\par

% structure: p0108 source=h3 target=subsection reason=relative-heading-rank; base=h2

\subsection{第51章 凡由亚当所生的人都在定罪之下，唯有在基督内重生的才能免于定罪。}

但稍后他说：\BibleQuote{如此，因一次的过犯，审判归于众人而定罪；同样，因一次的正义行为，恩赐归于众人而成义得生。}这足够清楚地显示，凡由\Person{亚当}所生的，无一不处于定罪之下，除非他在\Person{基督}内重生，无人能免于定罪。\par

% structure: p0110 source=h3 target=subsection reason=relative-heading-rank; base=h2

\subsection{第52章 在圣洗圣事中，即基督死亡与复活的象征，无论婴孩还是成年人，都死于罪恶，好在新生命中行走。}

在宗徒论及因一人而来的定罪和因一人而来的无偿恩赐，已按书信该部分所需充分阐述之后，他继续谈论在基督十字架上的圣洗圣事的伟大奥秘，并清楚地向我们阐明：在基督内的圣洗无非是基督之死的肖像，而基督在十字架上的死亡无非是罪之赦免的肖像；因此，祂的死如何真实，我们罪的赦免也同样真实；祂的复活如何真实，我们的成义也同样真实。他说：\BibleQuote{我们可说什么呢？我们岂可仍在罪中，好使恩宠增多吗？}因为他先前曾说：\BibleQuote{然而罪恶在哪里增多，恩宠在那里也格外增多。}因此，他自问：为着随之而来的格外增多之恩宠，是否应当继续留在罪中？但他回答：\BibleQuote{绝不！}并加上：\BibleQuote{我们这死于罪的人，岂能仍在罪中活着呢？}然后，为表明我们死于罪，他说：\BibleQuote{难道你们不知道，我们这受过洗归于耶稣基督的人，就是受洗归于祂的死吗？}于是，既然我们受洗归于基督的死证明我们死于罪，那么连受洗归于基督的婴儿也死于罪，因为他们受洗归于祂的死。因为没有任何例外：\BibleQuote{我们这受过洗归于耶稣基督的人，就是受洗归于祂的死。}说这话是为证明我们死于罪。那么，婴儿在\OriginalTerm{重生}{regeneration}中所死的，不是那出生时带来的罪，又是什么罪呢？因此，以下经文也适用于他们：\BibleQuote{所以，我们藉着洗礼已与祂同葬于死，为的是基督怎样藉着父的光荣从死者中复活，我们也怎样在新生活中度生。因为我们如果与祂同受相似于祂的死亡，也必与祂同受相似于祂的复活；我们知道，我们的\OriginalTerm{旧人}{old man}已与祂同钉十字架，为的是使罪恶的身体失效，使我们不再作罪的奴隶，因为已死的人是脱离了罪。所以，如果我们与基督同死，我们相信也必与祂同生；我们知道基督既从死者中复活，就不再死，死亡不再统治祂了。因为祂死，是死于罪一次；祂活，是活于天主。同样，你们也要看自己是在罪上死了，在耶稣基督内为天主而活。}他开始证明我们不可仍在罪中使恩宠增多时，曾说：\BibleQuote{我们这死于罪的人，岂能仍在罪中活着呢？}为表明我们死于罪，他加上：\BibleQuote{难道你们不知道，我们这受过洗归于耶稣基督的人，就是受洗归于祂的死吗？}他结束这段经文正如开始时一样。因为他引入基督的死，意在暗示基督自己也死于罪。祂死于什么罪呢？不是死于肉身吗？这肉身中没有罪，却有罪的肖像，因此被称为罪。所以，对那些受洗归于基督之死的人——这不只包括成年人，也包括婴儿——他说：\BibleQuote{同样，你们也要看自己是在罪上死了，在耶稣基督内为天主而活。}\par

% structure: p0112 source=h3 target=subsection reason=relative-heading-rank; base=h2

\subsection{第五十三章 基督的十字架与埋葬、复活、升天和坐在天主右边，是基督徒生活的肖像。}

因此，基督被钉十字架、被埋葬、第三日复活、升天、坐在圣父右边这一切事件，都是这样安排的，以致基督徒在此世的生活不仅是在奥秘意义上，而且是在实际上以它们为模范。因为关于祂的被钉十字架，经上说：\BibleQuote{凡属于基督的人，已把肉身连同私欲和偏情钉在十字架上了。}关于祂的埋葬：\BibleQuote{我们藉着洗礼与祂同葬于死。}关于祂的复活：\BibleQuote{为的是基督怎样藉着父的光荣从死者中复活，我们也怎样在新生活中度生。}关于祂的升天和坐在天父右边：\BibleQuote{你们既然与基督一同复活了，就该追求天上的事，那里有基督坐在天主的右边；你们要思念天上的事，不要思念地上的事。因为你们已经死了，你们的生命已与基督一同藏在天主内。}\par

% structure: p0114 source=h3 target=subsection reason=relative-heading-rank; base=h2

\subsection{第五十四章 基督的再来不属于过去，而是将在世界末日发生。}

但我们所信的关于基督将来行动的，即祂将从天上降临审判生者死者，这并不影响我们现在在此世的生活；因为那不是祂在地上所做的一部分，而是祂在世界末日要做的事。宗徒在上文所引经文之后立即提到这一点：\BibleQuote{当基督——我们的生命显现时，那时你们也要与祂一同在光荣中出现。}\par

% structure: p0116 source=h3 target=subsection reason=relative-heading-rank; base=h2

\subsection{第五十五章 \BibleQuote{基督要审判活人死人}这话可按两种方式理解。}

现在，\Quote{审判活人死人}这个说法可有两种解释：或者我们以\Quote{活人}指祂来临时尚未死亡、仍以肉身活着的人，以\Quote{死人}指那些已离世或将在祂来临前离世的人；或者我们以\Quote{活人}指义人，以\Quote{死人}指不义之人，因为义人和不义之人都将受审判。天主的审判有时在坏的意义上被理解，例如\Quote{行恶的复活定罪}；有时在好的意义上，例如\Quote{天主，求祢以祢的名拯救我，以祢的能力为我伸冤}。我们若想到天主的审判是将善人与恶人分开，并把善人安置在祂右边，使他们免受邪恶，不与恶人一同毁灭，便容易明白这一点。正因如此，圣咏作者呼求\Quote{天主，求祢为我伸冤}，随后又如同解释一般加上\Quote{和分辨我的案件与不虔敬的民族}。\par

% structure: p0118 source=h3 target=subsection reason=relative-heading-rank; base=h2

\subsection{第五十六章　圣神与教会。教会是天主的圣殿。}

现在，在适合信仰宣认的简要中论及了耶稣基督、天主的独生子、我们的主之后，我们接着说我们也信圣神——如此完成了构成神性的至高的天主圣三。然后我们提到圣教会。这样我们被引导理解：构成自由耶路撒冷的理性受造物，在言语次序上应隶属于创造者——至高的天主圣三；因为凡关于人而基督耶稣所说的，自然都与独生子的位格合一有关。因此，信经的真正次序要求教会隶属于圣三，如同房屋隶属于居住者，圣殿隶属于占据它的天主，城市隶属于建造者。我们这里应理解整个教会，不仅包括那在地上流徙的部分，从日出之地到日落之处赞美天主的圣名，并因从旧囚禁中获得释放而唱新歌；也包括那从受造之初就一直在天上忠于天主、从未经历堕落所带来的苦难的部分。这一部分由圣天使组成，他们享受不断的幸福；并且（正如当行的）他们帮助尚在流徙中的部分：因为这两部分将在永恒中成为一体，如今已在爱德中合一，整个教会受造为朝拜独一的天主。因此，整个教会或其任何部分，都不想被当作天主来朝拜，也不愿成为属于天主圣殿之人的天主——这圣殿是由那些被非受造的天主所创造的圣人所建成的。所以，圣神如果是受造物，就不能是创造者，而会是理性受造物的一部分。祂只是最高的受造物，因此就不会在信经中排在教会之前；因为祂自己会属于教会，属于天上那一部分。祂不会有圣殿，因为祂自己会是圣殿的一部分。但祂确有圣殿，正如宗徒所说：\BibleQuote{难道你们不知道，你们的身体是住在你们内的圣神的圣殿吗？你们这圣神是从天主得来的。}关于这身体，他在另一处说：\BibleQuote{难道你们不知道，你们的身体是基督的肢体吗？}那么，如果祂有圣殿，怎能不是天主呢？祂又怎能比基督小，而祂的肢体竟是基督呢？祂的圣殿并非与天主的圣殿不同，因为同一位宗徒说：\BibleQuote{难道你们不知道，你们是天主的圣殿吗？}并以此作证明：\BibleQuote{天主的圣神住在你们内。}因此，天主住在祂的圣殿内：不仅是圣神，还有圣父和圣子。圣子论及自己的身体——藉此祂成为在地上教会的元首（\BibleQuote{好使祂在万有中居首位}）——说：\BibleQuote{你们拆毁这座圣殿，三天之内我要把它重建起来。}所以，天主的圣殿，即整个至高的天主圣三的圣殿，就是圣教会，涵盖天地全境。\par

% structure: p0120 source=h3 target=subsection reason=relative-heading-rank; base=h2

\subsection{第五十七章　天上教会的状况。}

但对于那在天的教会部分，我们能说什么呢？只能说其中没有恶者，并且自从\BibleQuote{天主没有宽恕犯罪的天使，反而把他们投入地狱，交在黑暗的锁链中，以待审判}（如宗徒伯多祿所写），便无人从中堕落，也永不会有人堕落。\par

% structure: p0122 source=h3 target=subsection reason=relative-heading-rank; base=h2

\subsection{第五十八章　我们对天使团体的组织没有确切知识。}

至于那至福天上社会的组织如何，有何等级差异：尽管所有人都被\Quote{天使}这一总名称称呼——如我们在\BibleBook{}{希伯来书}中读到：\BibleQuote{天主曾向哪一个天使说过：\InnerQuote{你坐在我右边}？}（这措辞显然包括所有天使，无一例外）——我们却发现有些被称为\Quote{总领天使}；而总领天使是否与那称为\Quote{万军}的相同，以致\BibleQuote{你们所有天使，请赞美祂；你们所有万军，请赞美祂}这一说法等同于\Quote{你们所有天使，请赞美祂；你们所有总领天使，请赞美祂}？以及那四个名称——\BibleQuote{无论是座者，或是宰制者，或是率领者，或是掌权者}——宗徒似乎以此毫无例外地包括所有天上团体，它们各有什么涵义？让那些能回答这些问题的人回答吧，如果他们也能证明自己的答案正确；至于我，我承认无知。我甚至不能确定这一点：太阳、月亮和所有星辰是否属于同一团体，尽管许多人认为它们只是发光体，没有知觉或理智。\par

% structure: p0124 source=h3 target=subsection reason=relative-heading-rank; base=h2

\subsection{第五十九章。天使所取的身体引发一个非常困难且不太有用的讨论课题。}

此外，谁能说出天使以何等的身体向人显现，不仅使自己可见，甚至可触？\newline{}再者，他们如何不是通过物质的身体，而是借着属神的能力，将异象呈现给心灵的属灵之眼，而非肉身的眼睛；或者，如何将某些话语不是从外面送入耳中，而是从人灵魂的内在说出，而他们自己也停留在那里，正如先知所记载的：\BibleQuote{那在\Emph{我内}说话的天使对我说}（他并没有说\Quote{那对我说话的}，而是说\Quote{那在我内说话的}）；\newline{}或者，如何在睡眠中向人显现，并通过梦境传话，正如我们在福音中读到：\BibleQuote{看，上主的天使在梦中向他显现，说}？\newline{}因为这些沟通方式似乎暗示天使没有可触的身体，这就引出了一个非常难解的问题：圣祖们如何洗脚？雅各伯如何以如此确凿物质的方式与天使角力？\newline{}提出这样的问题，并尽己所能猜测答案，对心智是一种有益的锻炼，只要讨论保持在适当的范围内，并且避免误以为自己知道那些其实不知道的事。\newline{}因为，在这些问题及类似的问题上，我们有何必要去肯定、否定或精确界定呢？即使完全无知，也无可指责。\par

% structure: p0126 source=h3 target=subsection reason=relative-heading-rank; base=h2

\subsection{第六十章。当撒殚化身光明的天使时，能够识别他的诡计更为必要。}

当撒殚化身光明的天使时，我们更需要运用全部的分辨力和判断力，以免被他狡猾地引入有害的歧途。\newline{}因为，如果他只欺骗肉身的感官，而不扭曲那使人能过信德生活的真实而健全的心智，那么宗教便无危险；\newline{}或者，如果他假扮为善，行出与善天使相宜的事，说出与善天使相宜的话，而我们相信他是善的，这种错误对基督徒的信德并不有害或危险。\newline{}但是，当他通过那些与其本性相异的手段，进而将我们引向他自己的道路时，我们就必须格外警醒，识别并拒绝跟随他。\newline{}然而，除非天主约束并保护他们，有多少人能躲避他一切致命的诡计呢？\newline{}但正是这事的困难本身有益，因为它阻止人信赖自己或他人，引导众人唯独依靠天主。\newline{}无疑，任何虔诚的人都不会怀疑，这对我们最为有益。\par

% structure: p0128 source=h3 target=subsection reason=relative-heading-rank; base=h2

\subsection{第六十一章。地上的教会借着中保的血从罪中被救赎。}

那么，教会这一部分，即由圣天使和天主的众军组成的那部分，将在世界末日时，当我们与它一同分享永福时，才让我们认清其本来面目。\newline{}但另一部分，即与它分离、在世上漂泊作客的部分，我们更为了解，因为我们既属于它，又因它由人组成，而我们也是人。\newline{}教会的这一部分，已通过那位没有罪的中保的血，从一切罪中被救赎出来，它的歌声是：\BibleQuote{如果天主为我们，谁能反对我们呢？\newline{}祂没有怜惜自己的儿子，反而为众人把祂交出来。}\newline{}基督并非为天使而死。\newline{}然而，通过祂的死亡为人类所做的一切，在某种意义上也是为天使做的，因为罪在人神之间所设立的仇恨被消除了，友谊得以恢复，而通过人类的救赎，那伟大的叛变在天使行列中所留下的空缺也被填补了。\par

% structure: p0130 source=h3 target=subsection reason=relative-heading-rank; base=h2

\subsection{第六十二章。借着基督的祭献，万有得以恢复，天地之间得以和平。}

当然，圣天使受教于天主，他们正是在永恒默观祂的真理中得享幸福，知道有多少人类将补充他们的行列，填满他们全体公民的数目。\newline{}因此，宗徒说：\BibleQuote{万物都在基督内归于一体，无论是天上的，还是地上的。}\newline{}天上的事物被聚集，是指因天使的堕落而从天上失落的一切，从人类中得以恢复；\newline{}地上的事物被聚集，是指那些预定得永生的人，从旧日的败坏中被救赎出来。\newline{}如此，通过那唯一的一次祭献，即中保被献上的那次祭献——法律下众多祭牲所预表的唯一祭献——天上的事物与地上的事物和好，地上的事物与天上的事物和好。\newline{}因此，同一宗徒说：\BibleQuote{因为父乐意使一切的圆满居住在他内，\newline{}并借着他在十字架上所流的血立定了和平，\newline{}借着他使万有，无论是地上的还是天上的，都与自己和好。}\par

% structure: p0132 source=h3 target=subsection reason=relative-heading-rank; base=h2

\subsection{第六十三章。在天上掌权的天主的平安超越一切理解。}

此平安，正如圣经所说，\Quote{超越一切理解}，除非我们完全享有它，否则无法认识它。因为天上之物在何种意义上得以和好，除非它们与我们和好，\Emph{即}通过与我们和谐一致？因为在天空中，存在着的所有有智受造物之间，以及它们与造物主之间，有着不间断的平安。而此平安，如经上所记，超越一切理解；但这当然是指我们的理解，而非那些常仰瞻圣父面容者的理解。如今，无论我们人的理解力有多大，我们只是部分地认识，且透过镜子如同揣测。但当我们与天主的众天使平等时，我们就会像他们一样面对面地看见；我们将对他们怀有如同他们对我们那样大的平安，因为我们爱他们如同我们被他们所爱一样。因此，他们的平安将为我们所知：因为我们的平安将与他们相似，与他们一样大，那时也不会超越我们的理解。但天主的平安，即祂向我们怀有的平安，无疑不仅超越我们的理解，也超越他们的理解。而这是必然的：因为每个有福的理性受造物都从祂获得福乐；祂并不从受造物获得福乐。从这个观点看，最好将那段经文中的\Quote{一切}解释为\Quote{天主的平安超越一切理解}，意味着甚至不容许有例外，连圣天使的理解也不例外；唯一可以破例的是天主自己。因为，当然，祂的平安并不超越祂自己的理解。\par

% structure: p0134 source=h3 target=subsection reason=relative-heading-rank; base=h2

\subsection{第六十四章 罪过的赦免覆盖圣人的全部有死生命，这生命虽无罪债，却非无罪。}

但是，当我们的罪过获得赦免时，天使们此刻就与我们和好。因此，在信经的次序中，提及圣教会之后，便放置了罪过的赦免。因为地上的教会正是借此而站立；正是借此，那曾经失落又被寻回的，不再失落。因为，除去圣洗的恩宠，它是作为原罪的解毒剂而赐下的，使我们由出生所加给我们的，由重生得以解除（这恩宠也除去所有已在思想、言语和行为中所犯的本罪）：除去这一伟大恩惠——它起始于人的恢复，并在其中我们所有的罪债，无论是原罪还是本罪，都被洗去——我们余生从拥有理智之时起，就不断提供罪过赦免的机会，无论我们在义德上有多大进步。因为天主的子女，只要他们活在这死亡的身体中，就与死亡搏斗。尽管论到他们，确实说了：\Quote{凡被天主的圣神引导的，都是天主的子女}，然而他们被天主的圣神引导，但作为天主的子女迈向天主时，却有一个缺点：他们也被自己的精神引导，这精神被可朽坏的身体所拖累；并且，作为人的子女，在人情欲的影响下，他们退回到旧有的水准，因而犯罪。然而有区别。因为尽管每种大罪都是罪，但并非每种罪都是大罪。所以我们说，圣人们的生活，只要他们还在这个可死的身体中，可能找不到大罪；但正如宗徒若望所说：\Quote{如果我们说我们没有罪，我们就是欺骗自己，真理也不在我们内。}\par

% structure: p0136 source=h3 target=subsection reason=relative-heading-rank; base=h2

\subsection{第六十五章 天主赦免罪过，但以补赎为条件；补赎的特定时期已由教会的律法所规定。}

但是，即使是大罪本身，无论多大，在圣教会内也可以获得赦免；天主的仁慈，真正悔改的人绝不应失望，各人按自己罪过的程度。而在悔改的行为中，当所犯的大罪具有将罪人从基督的身体上切除的性质时，我们不应过多考虑时间的长短，而要考虑忧伤的程度；因为天主不轻视破碎和痛悔的心。但由于一个人的忧愁常常隐而不现，不通过言语或其他迹象为他人所知，当它向那位说\Quote{我的呻吟不会隐藏于你}的天主显明时，治理教会的人便合理地规定了补赎的时期，以使罪过得赦免的教会得到满足；而在教会之外，罪过不得赦免。因为唯有教会领受了圣神的保证，没有它就没有罪过的赦免——至少是那种使被赦免者获得永生的赦免。\par

% structure: p0138 source=h3 target=subsection reason=relative-heading-rank; base=h2

\subsection{第六十六章 罪过的赦免主要关乎将来的审判。}

如今，罪过的赦免主要关乎将来的审判。因为，就此生而言，圣经的这句话是真实的：\Quote{沉重的轭在亚当子孙身上，从他们离开母腹之日，直到他们回归万物之母之日。}因此我们看到，甚至婴儿在领受圣洗和重生之后，也遭受各种痛苦的折磨；这样我们就明白，在救恩圣事中所阐明的一切，所指的更是将来福乐的希望，而非保有或获得现世的祝福。因为许多罪过在此世似乎被忽视，未受惩罚，其惩罚被保留到将来（因为基督作为审判生者死者的日子特别被称为审判日，这并非徒然）；正如另一方面，许多罪过在此生受了惩罚，却仍被赦免，在后世不再带来惩罚。因此，关于某些现世的惩罚，即在此生加于罪人身上的，宗徒对那些罪过已被涂抹而不留于最后审判的人说：\Quote{因为我们若省察自己，就不会受审判；但我们受审判时，是受主的惩戒，免得我们与这世界一同被定罪。}\par

% structure: p0140 source=h3 target=subsection reason=relative-heading-rank; base=h2

\subsection{第六十七章 没有行为的信德是死的，不能使人得救。}

此外，有些人相信：那些不放弃\Person{基督}之名，在教会内受祂的洗礼，从未因任何分裂或异端被切断与教会联系的人，尽管他们生活在最严重的罪中，既不曾在补赎中洗去罪，也不曾以施舍赎罪，却固执于罪直到生命最后一天，他们将藉火得救；也就是说，虽然他们将遭受火的惩罚，时间长短与他们的罪行和恶行大小相称，但他们不会受永火的惩罚。但那些相信这一点、且身为天主教徒的人，在我看来，似乎被一种人性中固有的仁慈之情引入了歧途。因为当我们查阅圣经时，圣经给出了截然不同的答案。我曾就这个题目写过一本书，名为\WorkTitle{论信仰与工作}，在其中，尽我所能，在天主助佑下，我从圣经中指明：那拯救我们的信德，正是宗徒\Person{保禄}清楚描述的那种，他说：\BibleQuote{在\Person{耶稣基督}内，割损或未受割损都不算什么，唯有藉爱德行事的信德才算什么。}但如果它行恶而非行善，那么毫无疑问，正如宗徒\Person{雅各伯}所说，\BibleQuote{信德若没有行为，本身便是死的。}同一位宗徒又说：\BibleQuote{我的弟兄们，若有人说自己有信德，却没有行为，有什么益处？难道这信德能救他吗？}此外，如果仅凭信德，一个恶人将藉火得救，而这正是蒙福的宗徒\Person{保禄}所说的\BibleQuote{他自己可得救，但仍像从火中经过一样}，那么没有行为的信德就能救人，而他的同宗徒\Person{雅各伯}所说的就必定是假的。而且\Person{保禄}自己在另一处说的也必定是假的：\BibleQuote{不要自欺：无论是淫乱者、拜偶像者、奸淫者、作娈童者、好男色者、偷窃者、贪婪者、酗酒者、辱骂者、勒索者，都不能承继天主的国。}因为如果那些固执于这些恶行的人，竟能因他们对\Person{基督}的信德而得救，那么他们不能承继天主的国，这怎么可能是真的呢？\par

% structure: p0142 source=h3 target=subsection reason=relative-heading-rank; base=h2

\subsection{第六十八章. 关于\BibleRef{格前 3:11-15}段落中\Quote{自己可得救，但仍像从火中经过一样}的真意。}

但是，既然宗徒们这些最清晰明确的宣告不可能是假的，那么那段关于在根基（\Person{基督}）上建造，不是用金、银、宝石，而是用木、草、禾秸之人的晦涩说法（因为据说正是这些人将得救，但仍像从火中经过一样，根基的功劳拯救了他们），就必须如此解释，以便不与上面引用的清晰陈述相冲突。现在，木、草、禾秸可以无矛盾地理解为对世俗之事的依恋，尽管这些事物本身可能是合法的，但若失去它们，心灵就会忧伤。而且，虽然这种忧伤像火一样燃烧，但只要\Person{基督}在心灵中占据根基的位置——也就是说，没有任何事物被置于祂之上，而且这个人尽管被忧伤燃烧，却仍然更愿意失去他所深爱的事物，也不愿失去\Person{基督}——他便藉火得救。然而，如果在诱惑之时，他宁愿持守暂时和属世的事物而不愿持守\Person{基督}，他就没有以\Person{基督}为根基；因为他把属世的事物放在首位，而在建筑物中，没有什么是在根基之前的。再者，宗徒在此处所说的火，必须是两种人都要经历的火，即那在根基上建造金、银、宝石的人，以及那建造木、草、禾秸的人。因为他紧接着说：\BibleQuote{这火要试验各人的工程，看看是哪一种。若有人在根基上所建的工程存留，他必得赏报；若有人 的工程被烧毁，他就要受损失，但他自己可得救，但仍像从火中经过一样。}那么，这火所考验的，不仅是其中一人的工程，而是二人的工程。现在，逆境的考验是一种火，这在别处明确提到：\Quote{炉火试炼陶人的器皿，而逆境的炉火试炼义人。}而且，在今世的生活中，这火正是按宗徒所说的方式运作。如果它遇到两位信徒，一位\BibleQuote{挂虑主的事，怎样才能讨主的喜悦}，也就是说，在\Person{基督}这个根基上建造金、银、宝石；另一位\BibleQuote{挂虑世俗的事，怎样才能讨妻子的欢心}，也就是说，在同一个根基上建造木、草、禾秸——前者的工程不被烧毁，因为他没有将爱情倾注于那些失去后会引发忧伤的事物；但后者的工程被烧毁，因为那些以贪恋之心享受的事物，失去时不可能没有痛苦。但由于我们假设，即使是后者也宁愿失去这些事物而不愿失去\Person{基督}，并且他没有因害怕失去它们而离弃\Person{基督}，尽管当失去它们时他感到忧伤，他得救了，但仍像从火中经过一样；因为对他所爱并失去之物的忧伤燃烧着他。但这忧伤并没有颠覆他，也没有吞噬他；因为他被那不可动摇、不可朽坏的根基所保护。\par

% structure: p0144 source=h3 target=subsection reason=relative-heading-rank; base=h2

\subsection{第六十九章. 有些信徒将来经过炼狱之火并非不可能。}

并非不可能，类似之事甚至在今世之后也会发生。这是一个可以探究的问题，或可确定，或可存疑：是否某些信徒将经过某种炼火，并且根据他们曾以或多或少的虔诚依恋那会朽坏之物，将较快或较慢地被解脱出来。然而，这决不能适用于那些经上所说\BibleQuote{不能继承天国}的人，除非经过适当的悔改，他们的罪得赦免。当我说\Quote{适当的}，我意思是他们不可在施舍上不结果实；因为圣经如此强调这一德行，以致我们的主预先告诉我们，祂将不认可右边之人的任何功德，除非他们在此德上丰饶，也不指责左边之人的任何缺失，除非他们缺少此德，那时祂要对前者说：\BibleQuote{我父所祝福的，你们来，继承天国}，对后者说：\BibleQuote{可咒骂的，离开我，往永火里去。}\par

% structure: p0146 source=h3 target=subsection reason=relative-heading-rank; base=h2

\subsection{第70章。除非生命改变，否则施舍不能赎罪。}

然而，我们必须谨防有人以为重罪——就是那些不能继承天国的人所犯的——可以每天犯，每天用施舍来赎。生命必须改变得更好；施舍必须用来向天主赎过去的罪，而非为将来犯这类罪购买免罚。因为祂没有给任何人犯罪的特许，尽管祂在慈悲中可能会涂抹已经犯的罪，如果我们不忽略做出适当的补赎。\par

% structure: p0148 source=h3 target=subsection reason=relative-heading-rank; base=h2

\subsection{第71章。信徒的日常祈祷为日常玷污其生命的小罪做出补赎。}

如今，信徒的日常祈祷为那些暂时的、轻微的小罪做出补赎，这些罪是今生必然发生的。因为他可以说：\BibleQuote{我们的天父}，因为现在他由水和圣神重生，成为这样一位天父的子女。这祈祷确实除去日常生活中的极小罪。它也除去那些曾使信徒的生命非常邪恶，但现在他因悔改而改变，已经弃绝的罪，前提是他正如真实地说\BibleQuote{宽免我们的罪债}（因为总不缺需要被宽免的罪债），也同样真实地说\BibleQuote{如同我们宽免我们的负债者}；也就是说，前提是他言行一致：因为宽恕一个求宽恕的人，实在是施舍。\par

% structure: p0150 source=h3 target=subsection reason=relative-heading-rank; base=h2

\subsection{第72章。有许多种施舍，施舍它们有助于获得我们罪的赦免。}

根据这一解释原则，我们主的话\BibleQuote{你们把现有的施舍给人，看，一切都是洁净的}适用于人出于慈悲所做的一切有益行为。因此，不仅是那给饥饿者食物、给口渴者饮料、给赤身者衣服、给旅客住宿、给逃亡者庇护、探望病患和囚犯、赎回俘虏、帮助弱者、引导盲人、安慰忧伤者、医治病人、给迷途者指路、给困惑者建议、满足匮乏者需要的人——不仅这个人，而且那宽恕罪人的人也施舍；那用责打纠正，或用某种纪律约束他权力管辖下的人，同时从心里宽恕那伤害他的罪，或为之祈祷以求宽恕的人，也是施舍者，不仅因为他宽恕或为罪祈求宽恕，也因为他斥责和纠正罪人：因为在此，他也显示慈悲。如今，许多好处被加给不情愿的领受者，当考虑他们的益处而非快乐时；他们自己常常证明是自己的仇敌，而他们真正的朋友却是那些他们误认为仇敌，并在盲目中以怨报德的人。（的确，基督徒甚至不允许以恶报恶。）因此有许多种施舍，通过施舍它们，我们协助获得罪的赦免。\par

% structure: p0152 source=h3 target=subsection reason=relative-heading-rank; base=h2

\subsection{第73章。所有施舍中最大的，是宽免我们的负债者，并爱我们的仇敌。}

然而，没有哪一样比从心里宽恕那得罪我们的罪更大。因为对一个没有对你行恶的人心怀善意甚至行善，是相对微小的事。爱你的仇敌，总是对那愿你受损害、且若能行便加害于你的人心怀善意，并在有机会时行善，是更高的事，是至极善心的结果。这正是遵行天主的命令：\BibleQuote{爱你们的仇敌，善待那恨你们的，为那迫害你们的祈祷。}但是，鉴于这种心境只有天主的完美儿女才能达到，而且虽然每个信徒都应当追求它，并藉着向天主的祈祷和与自己认真的斗争，努力将自己的灵魂提高到这一标准，然而如此高尚的善度几乎不能属于这么一大群人，就是那些我们相信在使用\BibleQuote{宽免我们的罪债，如同我们宽免我们的负债者}这一祈祷时蒙垂听的人；考虑到这一切，无疑地，那隐含的承诺得以实现，如果一个人虽尚未达到爱仇敌的地步，但当被一个得罪他的人请求宽恕时，他从心里宽恕了他。因为当他祈祷\BibleQuote{如同我们宽免我们的负债者}时，他确实渴望自己被宽恕，也就是说，\BibleQuote{宽免我们的罪债}当我们乞求宽恕时，如同我们宽免我们的负债者当他们向我们乞求宽恕时。\par

% structure: p0154 source=h3 target=subsection reason=relative-heading-rank; base=h2

\subsection{第74章。天主不赦免那些不从心里宽恕他人之人的罪。}

现在，那个因自己的罪而被推动去向被他得罪的人请求宽恕的人，不可被视为一个敌人，以至于现在爱他如同他仍在积极敌对时那样困难。而那从心里不宽恕悔改并请求宽恕的罪人的人，不必以为他自己的罪被天主宽恕了。因为真理不能说谎。有哪位福音的读者或听者会没有注意到，那说过\BibleQuote{我是真理}的同一人也教导了我们这种祈祷的方式；并且为了将这特定祈求深印在我们心中，祂说：\BibleQuote{因为如果你们宽恕别人的过犯，你们的天父也必宽恕你们；但如果你们不宽恕别人的过犯，你们的父也必不宽恕你们的过犯}？这警告的雷声不能唤醒的人，不是睡着了，而是死了；然而那声音如此强大，甚至能唤醒死人。\par

% structure: p0156 source=h3 target=subsection reason=relative-heading-rank; base=h2

\subsection{第七十五章 邪恶与不信者若不重生，施舍也不能使他们洁净。}

确实，那些生活在严重邪恶中，又不关心改变生活和品行，却在自己所有罪行和恶习中不断施舍的人，他们从主的话中徒然自慰：\BibleQuote{你们把里面的施舍了，看哪，一切对你们都是洁净的。}因为他们不明白这话的含义。但为了让他们明白，让他们听听祂怎么说。因为我们在福音中读到如下的话：\BibleQuote{耶稣说话的时候，有一个法利塞人请祂吃饭；祂便进去坐席。那法利塞人看见，诧异祂饭前没有先洗手。主对他说：\InnerQuote{你们法利塞人把杯盘的外面洗得干干净净，但你们里面却充满劫夺与邪恶。无知的人哪！那造外面的，不是也造了里面吗？只要把里面的施舍了，看哪，一切对你们都是洁净的。}}我们是否要理解为，对于那些没有基督信仰的法利塞人，只要他们像这些人所认为的施舍那样施舍，即使他们从未信过基督，也未由水和圣神重生，一切就都是洁净的呢？但事实是，所有没有被基督的信仰洁净的人都是不洁净的，正如经上所说：\BibleQuote{以信德洁净了他们的心。}并且宗徒说：\BibleQuote{在污秽和不信的人，什么都洁净不了，连他们的理智和良心也都玷污了。}那么，如果法利塞人不信，即使他们施舍，怎能一切都洁净呢？他们又怎能成为信徒，如果他们不愿相信基督，并借祂的恩宠重生呢？然而他们听见的是真的：\BibleQuote{你们把里面的施舍了，看哪，一切对你们都是洁净的。}\par

% structure: p0158 source=h3 target=subsection reason=relative-heading-rank; base=h2

\subsection{第七十六章 要正确施舍，当从自己开始，先怜悯自己的灵魂。}

因为那愿意按正道施舍的人，应当从自己开始，先施舍给自己。因为施舍是怜悯的工作；最真实的话是：\BibleQuote{怜悯你的灵魂，使天主喜悦。}我们为此重生，好使我们能中悦天主，而天主正当地不喜悦我们从出生带来的东西。这是我们首先的施舍，我们给自己施舍，当我们因怜悯的天主的仁慈，发现自己可怜，并承认祂判决的公义——我们由此成为可怜，宗徒所说：\BibleQuote{判决是由一人而定罪}——又赞美祂大爱的伟大，同一恩宠的宣讲者说：\BibleQuote{天主向我们彰显祂的爱，当我们还是罪人时，基督就为我们死了}——这样，我们正确地判断自己的可怜，并以天主所赐的爱德爱天主，我们就过圣洁而有德行的生活。但法利塞人虽将一切出产的十分之一施舍，连最细微的也不放过，却忽视了正义和天主的爱，因此他们没有从家里开始施舍，首先没有怜悯自己。正是针对这种爱的秩序，经上说：\BibleQuote{爱你的邻人如同你自己。}所以，当我们主责备他们，因为他们使自己外面干净，里面却充满劫夺和邪恶时，祂建议他们，首先要运用每个人对自己所欠的爱德，清洁里面。祂说：\BibleQuote{但是，你们把里面的施舍了，看哪，一切对你们都是洁净的。}然后，为表明祂建议的是什么，而他们不愿做什么，以及祂没有忽视或忘记他们的施舍，祂说：\BibleQuote{但法利塞人有祸了！}好像祂要说：我确实建议你们施舍，以便使一切对你们洁净；\BibleQuote{但你们有祸了！因为你们把薄荷、芸香并各样菜蔬献上十分之一}；好像祂要说：我知道你们的这些施舍，你们不必以为我现在就这些事劝戒你们；\BibleQuote{却忽视了正义和天主的爱}——这施舍能使你们从一切内在的不洁中洁净，以至于连你们现在洗涤的身体对你们也成为洁净的。因为\Quote{一切}的意涵就是如此，包括内外的，正如我们在另一处读到：\BibleQuote{先清洁里面，好使外面也清洁。}但以免祂似乎轻视他们从地上果实所施的舍，祂说：\BibleQuote{这些是你们应当做的}——指正义和天主的爱——\BibleQuote{而不可放下那些}——指献上十分之一。\par

% structure: p0160 source=h3 target=subsection reason=relative-heading-rank; base=h2

\subsection{第七十七章 若愿施舍给自己，当逃避不义；因为喜爱不义的人恨恶自己的灵魂。}

那么，那些以为他们可以通过施舍——无论多么慷慨，无论在金钱或实物上——为自己购得特权，得以不受惩罚地继续其穷凶极恶的罪行和可憎的恶习的人，不必如此自欺。因为他们不仅犯这些罪，而且如此喜爱它们，以至于他们愿意永远继续犯下去，只要他们能不受惩罚地这样做。凡喜爱不义的，就是恨恶自己的灵魂；凡恨恶自己灵魂的，对他自己的灵魂不是仁慈，而是残忍。因为按照世界爱它，就是按照天主恨它。但如果他愿意向它施舍，这施舍会使一切为他洁净，他就会按照世界恨它，按照天主爱它。因为没有人施舍，除非他从不缺乏的那一位那里领受他所施舍的。因此经上说：\BibleQuote{他的仁慈必迎接我。}\par

% structure: p0162 source=h3 target=subsection reason=relative-heading-rank; base=h2

\subsection{第七十八章 什么是小罪什么是大罪，取决于天主的判断。}

现在，什么是小罪，什么是大罪，不是由人的判断来决定的，而是由天主的判断。因为很明显，宗徒们自己在某些罪上给予了宽许：例如，宗徒\Person{保禄}对已婚者说：\BibleQuote{你们不要彼此亏负，除非两相情愿，暂时分房，为专务祈祷和禁食；以后仍要同房，免得撒殚因你们的不能节制而诱惑你们。}现在，可能有人会认为，与配偶交合，不是为了生育子女（这是婚姻的大祝福），而是为了肉欲的快乐，并为了拯救不节制的人，使他们在软弱时不被引入大罪，即奸淫、通奸或其他羞于启齿的不洁的罪行（他们可能被欲望在撒殚的诱惑下引入），这些可能不被认为是罪。我说，如果不是宗徒接着加上：\BibleQuote{我说这话，是出于宽许，并不是出于命令}，这或许根本不会被看作罪。那么，当众承认只有靠宗徒的权威，才允许那些这样做的人时，谁能否认这是罪呢？另一个同类例子是他说的：\BibleQuote{你们中间有人与别人有了争讼，怎么竟敢在不义的人面前起诉，而不在圣者面前呢？}紧接着：\BibleQuote{然而你们若有了关于今世事情的争讼，就派教会中最轻微的人来裁判吗？我说这话，是为叫你们羞愧。难道你们中间竟没有一个有智慧的人，能在弟兄之间判断吗？可是，弟兄与弟兄互相控告，而且在不信主的人面前控告！}在这种情况下，可能有人会认为，与他人争吵不是罪，唯一的罪是希望到教会外去判决，要不是宗徒马上加上：\BibleQuote{所以，你们彼此控告，已经是全然的过失了。}并且，免得有人借口称自己有正当理由，正在受冤枉，只想让法官的判决消除他的冤屈，宗徒立即预见了这些想法和借口，说道：\BibleQuote{为什么你们不宁愿受屈？为什么你们不宁愿被亏待？}这样把我们带回到我们主的话：\BibleQuote{谁若愿意与你争讼，拿你的内衣，你连外衣也让他拿去。}以及：\BibleQuote{凡夺取你东西的，不要再索回。}因此，我们的主禁止祂的门徒为了世俗的事务与他人对簿公堂。秉承这一原则，宗徒在这里声明，这样做是\BibleQuote{全然的过失}。但是，尽管如此，他允许在教会内由其他弟兄裁判这样的弟兄间的案件，只严厉禁止带到教会外面去处理，显然这里再次向软弱的人延伸了宽许。因此，正是鉴于这些罪，以及其他同类罪，还有其他更微小的，包括言语和思想上的过犯（如宗徒\Person{雅各伯}所承认的：\BibleQuote{我们众人都犯许多过失}），我们需要每日常常向主祈祷，说：\BibleQuote{求你宽免我们的债}，并真诚地加上：\BibleQuote{如同我们宽免了欠我们债的人。}\par

% structure: p0164 source=h3 target=subsection reason=relative-heading-rank; base=h2

\subsection{第七十九章 看起来非常微小的罪，有时实际上非常严重。}

再者，有些罪若圣经不显示它们确实非常严重，就会被认为非常微小。因为若不是那真理本身这样说，谁会想到那称自己弟兄为\BibleQuote{愚妄的人}的人有地狱之火的危险呢？然而祂立即对伤口施以治疗，给出了与受冒犯的弟兄和好的规则：\BibleQuote{所以，你若在祭台前献你的礼物，在那里想起你的弟兄有什么怨你的事，就把礼物留在祭台前，先去与你的弟兄和好，然后再来献你的礼物。}再者，若不是我们从宗徒对这样的人所说的恐惧中评估出恶的严重性——他说：\BibleQuote{我怕我为你们白白地劳苦了}——谁会认为观察日子、月份、季节、年份是这么大的罪呢？如同那些因人的虚妄学说而认为某些日子、月份或年份是吉利或不吉利的，因而焦虑或不愿开始任何事情的人所做的那样。\par

% structure: p0166 source=h3 target=subsection reason=relative-heading-rank; base=h2

\subsection{第八十章 无论多么重大可憎的罪，当我们习惯于它们时，就会显得微不足道。}

此外，罪无论多么重大可憎，一旦人习以为常，就被视为轻微，甚至不以为是罪；以至于这样的罪不仅不被隐藏，反而被夸耀，四处张扬；正如经上所载：\BibleQuote{恶人因心里的欲望而自夸，贪婪的人赞颂自己，主却憎恶他。}圣经中称这种不义为\OriginalTerm{呼喊}{cry}。你在依撒意亚先知论恶葡萄园的例子中看到：\BibleQuote{祂期待公平，看，却是压迫；期待正义，看，却是呼喊。}因此创世纪中也有这话：\BibleQuote{索多玛和哈摩辣的呼喊声极大，}因为这些城邑中的罪行不仅不受惩罚，反而公然施行，仿佛受到法律保护。如今在我们这个时代，许多形式的罪，虽不完全与索多玛和哈摩辣相同，却已如此公开而习以为常，以至于我们不仅不敢将平信徒逐出教会，甚至不敢因这些罪而贬黜神职人员。以致几年前，当我讲解\WorkTitle{迦拉达书}，注释宗徒所说的那处\BibleQuote{我为你们害怕，唯恐我白白地为你们劳苦了}时，我不得不呼喊说：\Quote{祸哉，人的罪！因为只有在不习惯的时候，我们才躲避它们；一旦习以为常，尽管天主圣子的血曾倾流以洗净它们，尽管它们如此重大，完全关闭了天主的国，但持续的熟悉导致容忍一切，惯常的容忍导致行许多罪。主啊，求你使我们不至于行那些我们无力阻止的罪。}但我将看看，过度的悲伤是否使我言语轻率。\par

% structure: p0168 source=h3 target=subsection reason=relative-heading-rank; base=h2

\subsection{第八十一章 罪的两个原因：无知与软弱；我们需要天主的助佑来克服二者。}

现在我要说这一点，我曾在其他作品中多次说过。导致罪有两个原因：要么我们尚不知道自己的本分，要么我们知道了本分却不履行。前者是无知之罪，后者是软弱之罪。我们应当与这些争战；但是，除非得到天主的助佑，不仅使我们看清本分，而且在我们清楚看见时，使爱正义的心比爱世俗之物的心更强大——因为贪恋世俗之物或害怕失去它们，会让我们明知故犯——否则我们必定在战斗中失败。在后一种情况下，我们不仅是罪人（因为即使在无知中犯错，我们也是罪人），而且是违犯法律的人；因为我们不做知道该做的事，却做知道不该做的事。因此，我们不仅应当在犯罪后祈求宽恕，说：\BibleQuote{宽免我们的罪债，如同我们宽免得罪我们的人；}还应当祈求指引，使我们免受诱惑，说：\BibleQuote{不要让我们陷入诱惑。}我们要向那一位祈祷，正如圣咏作者所说的：\BibleQuote{上主是我的光明，我的救援：}祂是我的光明，因为祂消除我的无知；祂是我的救援，因为祂除去我的软弱。\par

% structure: p0170 source=h3 target=subsection reason=relative-heading-rank; base=h2

\subsection{第八十二章 天主的慈悲对于真正的悔改是必需的。}

甚至补赎本身，当按教会法律有充分理由进行时，也常因软弱而被规避；因为羞耻是对失去快乐的恐惧，当人的好评比使人谦卑自下做补赎的正义更能带来快乐时，人便会规避补赎。因此，天主的慈悲不仅在人悔改时需要，甚至在人悔改之前，也需要引领他悔改。否则如何解释宗徒论及某些人所说的话：\BibleQuote{或许天主会赐给他们悔改的心}？而在\Person{伯多禄}痛哭之前，福音记载：\BibleQuote{主转过身来，看了看\Person{伯多禄}。}\par

% structure: p0172 source=h3 target=subsection reason=relative-heading-rank; base=h2

\subsection{第八十三章 轻视天主慈悲的人犯了对圣神的罪。}

凡不相信教会中罪过得赦免，轻视天主这伟大慈悲的恩赐，并执拗至死的人，就犯了不得赦免的亵渎圣神的罪，因为圣神是基督借以赦罪的那一位。但我曾在一本专论此问题的书中，尽可能清楚地讨论了这一难题。\par

% structure: p0174 source=h3 target=subsection reason=relative-heading-rank; base=h2

\subsection{第八十四章 身体的复活引出了许多问题。}

现在，关于身体的复活——不是像有些人那样，暂时复活后又死去，而是像\Person{基督}自己的身体那样复活，进入永生——我不认为我能简短地讨论，同时又能对通常提出的所有问题给出满意的回答。然而，所有的人都将复活，无论是已生或将生的人，无论是已死或将死的人，这是每个基督徒都不应有丝毫怀疑的。\par

% structure: p0176 source=h3 target=subsection reason=relative-heading-rank; base=h2

\subsection{第八十五章 关于流产的胎儿。}

因此，首先出现了一个关于流产胎儿的问题，他们确实在母腹中诞生，但尚未达到可以重生（再生）的地步。因为如果我们断定这些胎儿也要复活，我们就不能反对任何关于已成形胎儿的推论。然而，谁不更倾向于认为未成形的流产胎儿就像从未结实的种子一样消亡呢？但是，谁敢否认——即使他不敢肯定——在复活时，一切形态上的缺陷都将被补全，时间本应带来的完美不会缺少，正如时间带来的瑕疵也不会存在？这样，本性既不会缺少与它相合的长久岁月所应添加的东西，也不会因长久岁月所添加的相反之物而受损；未完成的将被完成，正如受损的将被更新。\par

% structure: p0178 source=h3 target=subsection reason=relative-heading-rank; base=h2

\subsection{第八十六章 如果他们曾经活着，当然也必已死去，因此他们在死人复活中必有分。}

因此，下列问题可由博学之士非常仔细地查考和讨论，尽管我不知道人是否有能力解决它：婴儿何时开始在母胎中存活？生命是否在活物动作显明之前就以潜伏的形式存在？否认那些为了不让母亲因胎死腹中而丧命而被从子宫中肢解取出的胎儿曾经活过，似乎太过大胆。一个人从开始活着的时刻起，他就可能死亡。若他死了，无论死亡在何处降临到他，我找不到任何理由可以否认他有权分享死人的复活。\par

% structure: p0180 source=h3 target=subsection reason=relative-heading-rank; base=h2

\subsection{第八十七章 论畸形生产之案例}

我们甚至不可断言，那些生而存活、尽管很快死去的畸形儿，不会复活，或会带着他们的畸形复活，而非带着修正和完善的身体复活。但愿天主禁止！最近东方生了一个双肢人，有最可信赖的见过他的弟兄带来报告——真福司铎\Person{热罗尼莫}曾留下书面记录——但愿天主禁止我们以为在复活时会有一个人长着双肢，而非两个不同的人，如同生了双胞胎那样。因此，其他生产，或因多余、或缺损、或严重变形而被称作\Emph{畸形}的，在复活时都将恢复人的正常形态；每一个单独的灵魂将拥有自己的身体；即使身体原先是连在一起的，也不会连在一起，而是各自单独拥有构成完整人身体的所有肢体。\par

% structure: p0182 source=h3 target=subsection reason=relative-heading-rank; base=h2

\subsection{第八十八章 论身体的质料永不消失}

构成人可朽身体的尘世质料也永不消失；尽管它可能化为尘土和灰烬，或分解为水汽和蒸气，或转变为其他物体的实体，或分散到元素中，或成为野兽或人的食物并转化为它们的肉，但在瞬间它都会回归到最初赋予它生命、使之成为人并生长的那一个人灵魂。\par

% structure: p0184 source=h3 target=subsection reason=relative-heading-rank; base=h2

\subsection{第八十九章 但此质料在复活的身体中可有不同安排}

这尘世质料，当灵魂离开后成为尸体，在复活时不会这样恢复，以至于它分离成各部分、并以各种形式和外观成为其他事物的一部分（尽管它们都将回到它们曾分离出去的同一身体），却必须回到它们原本所在的身体同一部位。否则，设想头发恢复我们频繁剪剃所失去的一切，指甲恢复我们常修剪的一切，就会给人想象出一幅如此丑陋和畸形的画面，以至于使身体的复活几乎不可信。正如一尊由可熔金属铸造的雕像，若被火烧熔、击碎或变成一团不成形的东西，而一位雕塑家想用同一数量的金属修复它，那么对于作品的完整性来说，金属的任何微粒被放在雕像的哪个部位都无关紧要，只要修复的雕像包含了原雕像的所有材料；同样，天主，那位奇妙而不可言喻能力的工匠，将以奇妙而不可言喻的速度恢复我们的身体，使用它原本构成的全部材料。无论头发是否回到头发，指甲是否回到指甲，还是这些已经消逝的部分变为肉，并被召唤到身体的其他部位，这都不会影响其恢复的完整性，因为这位伟大的工匠谨慎地确保没有任何东西不适当或不得其所。\par

% structure: p0186 source=h3 target=subsection reason=relative-heading-rank; base=h2

\subsection{第九十章 若复活者身体之间存在差异和不平等，亦无任何冒犯或不协调}

那些复活的人不一定因为生前身高不同而具有不同的身高；也不一定瘦子复活时保持原先的瘦，胖子保持原先的胖。但若造物主的设计是各人保留自己面容的特性并保持与生前可辨认的相似，而在其他身体优势上人人平等，那么各人所构成的质料可以被修改，以致没有失落任何部分，任何缺陷都由那能凭空创造祂所愿者来补足。但若复活者身体中存在井井有条的不平等，如同构成完整和声的各个声音，那么各人身体的质料将被如此处理，以形成一个适合天使集会的人，并且不会在其中带来任何使人感官不快的东西。确实，那里不会有任何不体面的事；无论那里有什么，都将优雅而合宜：因为若有不体面的事，它就不会存在。\par

% structure: p0188 source=h3 target=subsection reason=relative-heading-rank; base=h2

\subsection{第九十一章 论圣人们复活后的身体将是属神的身体}

圣人们的身体将复活，摆脱一切缺陷、一切瑕疵，以及所有朽坏、沉重和阻碍。因为他们的行动将如他们的幸福一样完全。因此，他们的身体被称为\Emph{属神的}，尽管它们无疑将是身体，而非神灵。因为正如现在身体被称为\Emph{有灵的}，虽然它是身体，而非\OriginalTerm{灵魂}{anima}，同样，那时身体将被称为属神的，虽然它是身体，而非神灵。因此，就那如今压着灵魂的朽坏，以及那激动肉体去与精神争战的恶习而言，那时它不再是肉体，而是身体；因为有称为天上的身体。所以经上说：\BibleQuote{血肉不能继承天主的国}；并且，好似解释这一点，又说：\BibleQuote{朽坏不能继承不朽坏}。宗徒先称为\BibleQuote{血肉}的，后来称为\BibleQuote{朽坏}；先称为\BibleQuote{天主的国}的，后来称为\BibleQuote{不朽坏}。但就实体而言，那时它仍然是肉体。因为甚至在复活后，基督的身体仍被称为肉体。然而宗徒说：\BibleQuote{所种的是属生灵的身体，复活的是属神的身体}；因为那时肉体与精神之间的和谐将如此完美，精神无需任何滋养即可使臣服的肉体存活，以至于我们本性的任何部分都不会彼此分歧；而是，正如我们将摆脱外在的仇敌，对内我们也无己为敌。\par

% structure: p0190 source=h3 target=subsection reason=relative-heading-rank; base=h2

\subsection{第九十二章 恶人的复活。}

但至于那些从由第一人的罪所导致的毁灭群体中，未被那介于天主与人之间的唯一中保所救赎的人，他们也将复活，每人带着自己的身体，但只是为了与魔鬼及其使者一同受罚。至于他们将带着自己所有的疾病和身体畸形复活，带着他们在此世所拥有的病态与残缺的肢体，探问这些是徒劳的。因为当他们的永罚已成定局，我们无需耗费精力去揣测他们的健康或美丽，这些是不确定的事。也无需探究，如果他们的身体能感受痛苦，如何算作不朽坏的？或者如果免于死亡的可能，如何算作可朽坏的？因为没有真正的生命，除非有生命中的幸福；没有真正的不朽坏，除非健康不被任何痛苦破坏。然而，当不幸者不被允许死去时，那么，容我这样说，死亡本身也就不死；当痛苦不间断地折磨灵魂，永无止境时，朽坏本身也未完成。这在圣经中称为\BibleQuote{第二次死亡}。\par

% structure: p0192 source=h3 target=subsection reason=relative-heading-rank; base=h2

\subsection{第九十三章 第一次死亡和第二次死亡都是罪恶的后果。惩罚与罪过成比例。}

无论是第一次死亡——即灵魂被迫离开身体时发生，还是第二次死亡——即灵魂不被允许离开受苦的身体时发生，若无人犯罪，都不会加给人。当然，最轻微的痛苦将落在那些没有在带来的原罪上增加任何本罪的人身上；至于其余增加了本罪的人，每个人的惩罚在来世将更为可忍受，按其今世的不义较小而定。\par

% structure: p0194 source=h3 target=subsection reason=relative-heading-rank; base=h2

\subsection{第九十四章 圣人们在来世将更充分认识因恩宠所领受的恩惠。}

如此，当被弃绝的天使和人受永罚时，圣人们将更充分认识因恩宠所领受的恩惠。然后，在默观实际情况时，他们将更清楚地看到圣咏中那句话的意思：\BibleQuote{我要歌唱仁慈和公道}；因为唯有不配得的仁慈才使人得救赎，唯有应得的审判才使人被定罪。\par

% structure: p0196 source=h3 target=subsection reason=relative-heading-rank; base=h2

\subsection{第九十五章 天主的审判那时将得到解释。}

那时，许多如今模糊的事将变得清楚。例如，当两个婴儿，他们的情形在所有方面都看似相同，一个因天主的仁慈被拣选归己，另一个因祂的公义被舍弃（其中被拣选的那个可以认识到，若非仁慈干预，公义本该加于己身）；为何这两个之中，这一个被拣选而非那一个，对我们来说是一个无法解决的问题。又比如，为何奇迹没有行在那些看到奇迹就会悔改的人面前，却行在那些明知不会悔改的人面前？因为我们的主最明确地说：\BibleQuote{祸哉，你，\Place{苛辣匝因}！祸哉，你，\Place{贝特赛达}！因为在你们中间所行的异能，若行在\Place{提洛}和\Place{漆冬}，他们早已披麻蒙灰悔改了。}而且，天主不愿意他们得救，这决无不义，尽管若祂愿意，他们本可得救。那时，将以最清晰的智慧之光看到，如今对虔诚人来说是一种信德，虽然尚未成为确切的知识：天主的旨意是多么确定、多么不变、多么有效；祂能做多少祂不愿意做的事，但祂愿意的没有不能做的；圣咏作者的歌是多么真实：\BibleQuote{但我们的天主在天上，祂随己意行事。}如果天主曾经愿意某事却没有实现，这话当然不真；更糟的是，如果人的意志阻碍了全能者做祂所愿的。因此，除非全能者的旨意（祂或允许它发生，或亲自做它），没有任何事发生。\par

% structure: p0198 source=h3 target=subsection reason=relative-heading-rank; base=h2

\subsection{第九十六章 全能的天主即使在允许恶事上也做得很好。}

我们也不能怀疑，即使是在容许恶的存在这件事上，天主也做得很好。因为祂只是在祂审判的正义中容许它。而凡是正义的，必定是善的。因此，虽然恶就其本身而言并非善，但恶与善并存这件事本身却是善。因为如果恶的存在不是一件好事，那么全能之善的天主就不会容许它存在；而祂既能轻易地拒绝容许祂所不愿的事，也能成就祂所愿的事，这一点是毫无疑问的。如果我们不信这一点，我们信经的第一句话就危险了，因为我们在那里宣认相信全能者天主圣父。因为如果祂不能成就祂所乐意的一切，或者祂的全能旨意被任何受造物的意志所阻碍，祂就不能真正被称为全能的。\par

% structure: p0200 source=h3 target=subsection reason=relative-heading-rank; base=h2

\subsection{第97章 宗徒说\BibleQuote{天主愿意所有的人都得救}，而事实上并非所有人都得救，这是在什么意义上说的？}

因此，我们必须探究，宗徒极其真实地论及天主的话是在什么意义上说的：\BibleQuote{祂愿意所有的人都得救}。因为事实上，并非所有人，甚至大多数人，都没有得救；这样看来，似乎天主的旨意没有成就，人的意志干预并阻碍了天主的旨意。当我们问为什么不是所有人都得救时，通常的回答是：\Quote{因为人自己不愿意。}这当然不能用在婴儿身上，因为愿意或不愿意不在他们的能力之内。但如果我们能将在受洗时他们所做的幼稚动作归因于他们的意志——当他们尽其所能地反抗时——我们就得说，就连他们也不愿意得救。然而，我们的主在福音中明确地说，当祂谴责那不信的城市时：\BibleQuote{我多少次愿意聚集你的子女，有如母鸡把自己的雏鸡聚集在翅膀下，而你们却不愿意！}似乎天主的旨意被人的意志战胜了，当最软弱的人以其不愿意而挡在路上时，最强大者的旨意就无法实现。那么，那位在天上和地上随意行事的全能者在哪里呢？如果天主愿意聚集耶路撒冷的子女，却没有成就？或者更确切地说，是耶路撒冷不愿意她的子女被聚集？但即使她不愿意，祂仍然聚集了祂所愿意的那么多子女：因为祂并不是有些事愿意并去做，另一些事愿意却不去做；而是\BibleQuote{祂在天上和地上，凡祂喜欢的，都一一实行}。\par

% structure: p0202 source=h3 target=subsection reason=relative-heading-rank; base=h2

\subsection{第98章 预定得永生完全出于天主的自由恩宠。}

再者，谁会如此愚蠢和亵渎，竟敢说天主不能按照祂的意願、在任何时候、任何地方，改变人的邪恶意志，使之转向善呢？但当祂这样做时，是出于仁慈；当祂不这样做时，是出于公义——因为\BibleQuote{祂愿意恩待谁，就恩待谁；祂愿意使谁硬心，就使谁硬心}。当宗徒说这话时，他是在阐明天主的恩宠，关于这一点，他刚刚提到了黎贝加腹中的双胞胎：\BibleQuote{他们还没有出生，还没有行善或作恶，为使天主选拔人的计划坚定不移，不因人的行为，而因召叫人的天主，就对她说：'年长的要服侍年幼的。'}关于此事，他又引用了另一段先知的话：\BibleQuote{我爱了雅各伯，而恨了厄撒乌。}但他意识到他所说的话可能会影响那些无法凭借悟性深入理解这恩宠之奥妙的人，于是他说道：\BibleQuote{那么，我们可说什麼呢？难道天主不公道吗？绝对不是！}因为在没有任何善恶行为的功过时，天主爱一个、恨另一个，似乎不公道。如果宗徒希望我们理解为一方有未来的善行、另一方有未来的恶行——这当然是天主预知的——他就不会说\BibleQuote{不因行为}，而会说\BibleQuote{不因未来的行为}，这样困难就解决了，或者更确切地说，根本就不会有困难。然而，在回答了\BibleQuote{绝对不是}之后——即天主绝对不公道——他接着证明天主这样做并没有不公义，并说：\BibleQuote{因为祂对梅瑟说：'我要恩待谁，就恩待谁；我要怜悯谁，就怜悯谁。'}现在，除了傻瓜，谁会认为天主不公义呢？无论祂是对那些应得惩罚的人施行惩罚的公义，还是对那些不配的人施以仁慈。然后他得出结论：\BibleQuote{这样看来，这不在于人愿意，也不在于人努力，而在于天主的仁慈。}所以那双胞胎生来都是义怒之子，并非因为他们自己的任何行为，而是因为他们被束缚在来自亚当的原罪处罚的锁链中。但那位说\BibleQuote{我要恩待谁，就恩待谁}的天主，出于白白的恩宠爱了雅各伯，出于应得的审判恨了厄撒乌。因为这审判是两人都该得的，前者从后者的例子中认识到，同样的惩罚没有降临到自己身上，并不能让他有理由因自己的任何功绩而夸耀，而只能夸耀于天主的恩宠之丰富；因为\BibleQuote{这不在于人愿意，也不在于人努力，而在于天主的仁慈}。的确，圣经的全部面貌，甚至可以说，其面容的每一线条，都通过极其深奥的类比，向每一个仔细研读的人传达着这个有益的警告：\BibleQuote{凡要夸耀的，应因主而夸耀}。\par

% structure: p0204 source=h3 target=subsection reason=relative-heading-rank; base=h2

\subsection{第99章 正如天主的仁慈是白白的，祂的审判也是公义的，不容置辩。}

在赞扬天主的仁慈之后，他说：\BibleQuote{如此，不是由于那愿意的，也不是由于那奔跑的，而是由于那显示仁慈的天主。}为了也赞扬祂的正义（因为那未获得仁慈的人所遇到的不是不义，而是正义，因为天主没有不义），他立刻补充说：\BibleQuote{因为圣经对法老说：\InnerQuote{我将你兴起，正是为这个缘故，为要在你身上显示我的权能，并使我的名传遍全地。}}然后他得出适用于两者的结论，即祂的仁慈和祂的正义：\BibleQuote{因此，祂对谁愿意施予仁慈就施予仁慈，祂愿意使谁心硬就使谁心硬。}\BibleQuote{祂施予仁慈}是由于祂的大善，\BibleQuote{祂使心硬}则毫无不义；如此，那被宽恕的人不能夸耀自己的功德，那被定罪的人也只能抱怨自己的缺失。因为正是恩宠单独将得救者与丧亡者分开，所有人都因共同根源而陷入同一普遍的丧亡。现在，如果有人听到这些，就说：\BibleQuote{祂为什么还指责呢？谁能抗拒祂的旨意呢？}好像人不应当因作恶而受责备，因为天主对谁愿意施予仁慈就施予仁慈，使谁心硬就使谁心硬。我们绝不应羞于用如同宗徒所回答的来回答：\BibleQuote{你这个人哪！你是谁，竟敢向天主顶嘴？受造物岂能对造物主说：\InnerQuote{你为什么这样造了我？}难道陶工没有权柄用同一团泥，造一个贵重的器皿，又造一个卑贱的器皿吗？}现在有些愚蠢的人认为宗徒在此没有给出回答；由于无法给出理由，他就斥责了提问者的僭越。但是这句话有极大的分量：\BibleQuote{你这个人哪！你是谁？}在这种事情上，它用一个词提醒人他的能力的极限，同时实际上传达了一个重要的理由。因为如果一个人不明白这些事，他怎能与天主争辩呢？如果他明白了，他就不再有什么可争辩的。因为那时他领悟到，整个人类在其叛逆的头颅中受到了如此正义的神圣审判，以至于如果人类中没有一个人得救，没有人能公正地质疑天主的正义；而且，那些得救的人应当以这样的方式得救，即通过更多未得救且留在其正义定罪中的人，显示整个人类应得的下场，以及天主的公正审判将把得救者引向何方，如果不是祂不应得的仁慈介入的话，以致所有想要夸耀自己功德的口都被堵住，而那夸耀的，应当夸耀于主。\par

% structure: p0206 source=h3 target=subsection reason=relative-heading-rank; base=h2

\subsection{第100章 天主的旨意从不被挫败，尽管许多行事与祂的旨意相反。}

这些是主的伟大工程，依照祂的一切喜悦而被寻求，并且如此明智地被寻求，以致当理性的受造物，无论是天使还是人类，犯罪时，不是遵行祂的旨意而是自己的旨意，祂竟利用那与造物主旨意相悖的受造物的意志，作为执行祂旨意的工具，那至善的天主就这样甚至将邪恶转为善用，为了定罪那些祂按正义预定受罚的人，并为了拯救那些祂按仁慈预定得恩宠的人。因为就受造物自身的意识而言，他们做了天主不愿发生的事；但就天主的全能而言，他们绝对无法实现他们的目的。因为正是在他们行事与祂的旨意相悖的事实中，祂关于他们的旨意反而得以实现。因此，\BibleQuote{主的伟大工程，依照祂的一切喜悦而被寻求}，因为以一种无法言说的奇妙方式，甚至那与祂的旨意相悖的事也不能挫败祂的旨意。因为如果不经祂许可（当然祂的许可不是不情愿，而是甘愿），这事就不会发生；并且一位善的天主不会允许邪恶发生，除非祂能在祂的全能中将邪恶转为善。\par

% structure: p0208 source=h3 target=subsection reason=relative-heading-rank; base=h2

\subsection{第101章 天主的旨意总是善的，有时通过人的邪恶意志得以实现。}

然而，有时候，一个人怀着良善的意志渴望天主不希望的事，尽管天主的意志也是善的，甚至更充分、更确定地是善的（因为祂的意志绝不会是邪恶的）：例如，如果一个好儿子焦虑地希望他的父亲活着，而天主的善意却是他应该死去。反之，一个人可能怀着邪恶的意志渴望天主因祂的良善所愿的事：例如，如果一个坏儿子希望他的父亲死去，而这也正是天主的意愿。显然，前者希望的事是天主所不愿的，而后者希望的事是天主所愿的；然而，前者的孝爱之情虽与天主的意愿不同，却比后者的不孝之情更符合天主的善愿，尽管后者的意愿与天主的相同。因此，在判断一个人的愿望是应受赞许还是反对时，必须考虑什么是适合人渴望的，什么是适合天主渴望的，以及每种情况下意志的真正动机。因为天主通过恶人的邪恶欲望完成祂的一些计划（当然是全部善好的）：例如，基督被杀害是通过犹太人的邪恶计谋，实现了圣父的善好旨意，而这事是如此善好，以至于当宗徒伯多禄表示不愿它发生时，那位将要被杀的称他为撒殚。那些虔诚信徒的好意看起来多么美好，他们不愿保禄上耶路撒冷，以免阿加波所预言的那恶事在那里发生在他身上！然而，天主的旨意是他应该为了宣讲基督的信仰而遭受这些恶事，从而成为基督的见证人。而祂的这个善好的旨意，并没有通过基督徒的善意谋划来实现，而是通过犹太人的恶谋，以至于那些反对祂旨意的人比那些自愿作为工具完成它的人更事奉祂。\par

% structure: p0210 source=h3 target=subsection reason=relative-heading-rank; base=h2

\subsection{第102章 全能天主的旨意绝不会挫败，也绝不会邪恶。}

然而，无论天使或人的意图多么强烈，无论善恶，无论这些意图符合天主的旨意还是与它相悖，全能者的旨意绝不会挫败；祂的旨意绝不可能是邪恶的；因为即使它施加邪恶，它也是公义的，而公义的当然不是邪恶的。因此，全能的天主，无论是在怜悯中祂愿意怜悯谁，还是在审判中祂使谁硬心，祂所行的绝不会不义，祂所作的无不出于祂自己的自由意志，而且祂从未愿意任何祂不实行的事。\par

% structure: p0212 source=h3 target=subsection reason=relative-heading-rank; base=h2

\subsection{第103章 对\BibleBook{弟茂德}{前书}2:4\BibleQuote{祂愿意所有的人都得救}这句话的解释。}

因此，当我们听到并读到圣经上说祂\BibleQuote{愿意所有的人都得救}时，虽然我们深知并非所有人都得救，但我们不应因此限制天主的全能，而应这样理解圣经：\BibleQuote{祂愿意所有的人都得救}的意思是，除非天主愿意某人得救，否则无人得救；并非祂不愿某些人得救，而是无人能脱离祂的旨意而得救；因此，我们应当祈求祂愿意我们得救，因为如果祂愿意，就必定实现。宗徒正是在论及向天主的祈祷时使用了这句话。同样，我们按同一原则解释福音中的话：\BibleQuote{那照亮每人的真光，正在进入这世界}；这不是说没有人不受光照，而是说没有人能不受祂光照。或者，经文说\BibleQuote{祂愿意所有的人都得救}，不是指祂不愿某些人得救（否则如何解释祂不愿在有些人面前行奇迹，而祂曾说他们若见了奇迹便会悔改呢？），而是我们应当将\BibleQuote{所有的人}理解为人类中所有阶层和境况的多样性——君王与臣民、贵族与平民、高者与低者、学者与愚者、健全者与虚弱者、聪明者与迟钝者、愚蠢者、富者与穷者、中产者；男性与女性、婴孩与孩童、青年；少年、中年与老年；每种语言、每种风尚、所有技艺、所有职业，连同意志与良心的一切无数差异，以及任何其他使人区分的事物。因为在所有这些类别中，难道有天主不愿其得救的人吗？祂通过其独生子、我们的主，愿意万民得救，并且确实拯救了他们；因为全能者无法白愿，凡祂所愿的必成就。宗徒曾嘱咐为众人祈祷，并特别加上\BibleQuote{为君王和一切有权位的人}，这些人可能因世俗地位的骄傲与浮华而回避基督徒信仰的谦卑。接着他说：\BibleQuote{这原是美好的，并在我们的救主天主面前是蒙悦纳的}，即应为这样的人祈祷；他立刻加上，似乎为消除任何绝望的借口：\BibleQuote{祂愿意所有的人都得救，并得以认识真理。}因此，天主以其伟大的屈尊俯就，认为让谦卑者的祈祷促成尊高者的得救是好的；我们确实有许多这样的例子。我们的主也在福音中使用了同样的表达方式，当祂对法利塞人说：\BibleQuote{你们把薄荷、芸香和各样蔬菜献上十分之一}。因为法利塞人不献他人的财产，也不献所有别地居民的所有蔬菜。因此，正如在这里我们必须将\BibleQuote{各样蔬菜}理解为各种蔬菜，同样在前一处经文里，我们可将\BibleQuote{所有的人}理解为各类人。我们也可用任何其他方式解释它，只要不被迫相信全能的天主曾愿意做某事而没有做成：因为撇开一切歧义，如果\BibleQuote{祂在天上地下都随己意而行}，正如圣咏作者颂唱的那样，祂当然没有愿意做任何没有做成的事。\par

% structure: p0214 source=h3 target=subsection reason=relative-heading-rank; base=h2

\subsection{第104章 天主预知第一个人的罪，从而安排了祂自己的计划。}

因此，如果天主预见到人会保持他受造时的无辜状态中的坚定意志，祂本愿意甚至将第一个人保存在他受造时的救恩状态中，并在他有生子之后，在适当时机免于死亡而迁移到一个更好的地方，在那里他不仅没有罪，甚至没有犯罪的欲望。但由于祂预见到人会滥用他的自由意志，即会犯罪，天主安排了自己的计划，更着眼于在人的罪恶中仍对他施恩，从而使全能者的善意不致因人的恶意而落空，反而能克服它得以实现。\par

% structure: p0216 source=h3 target=subsection reason=relative-heading-rank; base=h2

\subsection{第105章 人被创造得如此：既能择善，也能择恶；但在来世，选择恶将不可能。}

人最初被造时，适宜于既能愿意正确的事，也能愿意错误的事；若愿意前者，并非没有赏报；若愿意后者，并非没有惩罚。但在来世，他将不能愿意作恶；然而这并不限制他意志的自由。相反地，当他完全不可能成为罪的奴隶时，他的意志将更加自由。我们决不应责怪意志，或说它不再是意志，或不再被称为自由的，当我们如此渴望幸福时，不仅我们躲避痛苦，而且发现完全不可能做别的事。那么，正如灵魂现在不可能渴望不幸，同样将来它完全不可能渴望犯罪。但是，天主的安排不可破坏，按照这安排，他愿意显示一个甚至能够避免犯罪的理性受造物是多么好，而一个完全不能犯罪的受造物更是多么好；正如那是一种低级的不朽，却仍是不朽，当时人有可能避免死亡，尽管将来还有一个更完善的不朽，那时人将不可能死亡。\par

% structure: p0218 source=h3 target=subsection reason=relative-heading-rank; base=h2

\subsection{106、天主的恩宠在人类堕落前后均为救恩所必需。}

前者（即第一种不朽）人因运用自由意志而丧失；后者他将藉恩宠获得，而如果他未曾犯罪，他本可因功德获得。然而即使在那情况下，没有恩宠也不可能有功德；因为虽然仅仅运用人的自由意志足以招致罪恶，但人的自由意志不足以使他坚守义德，除非天主藉着赐予一部分他不改变的良善来帮助它。正如人有能力随时死亡（姑且不谈其他方式，任何人仅靠绝食就可结束生命），但若没有食物和其他生活资料，单纯的意志不能维持生命；同样，在乐园里，人能够仅凭其意志，单靠放弃义德来毁灭自己；但要维持义德的生活，单凭他的意志是不够的，除非有造物主的力量支持。然而，在堕落之后，需要更丰盛的天主仁慈，因为意志本身必须从罪恶和死亡的束缚中解放出来。而意志的自由丝毫不是来自自身，唯独来自因信耶稣基督而来的天主的恩宠；因此，我们藉以接受引导我们走向天主永恒恩赐的一切其他恩赐的意志本身，乃是由主所预备的，正如圣经所说。\par

% structure: p0220 source=h3 target=subsection reason=relative-heading-rank; base=h2

\subsection{107、永生虽为善工之赏报，其本身却是天主的恩赐。}

因此，就连永生本身——它确实是善工的赏报——宗徒也称之为天主的恩赐。他说：\BibleQuote{因为罪的工资是死亡}，\BibleQuote{但天主的恩赐是在我们的主基督耶稣内的永生。} \OriginalTerm{工资}{stipendium}是作为军役的报酬支付的；它不是恩赐：因此他说，\BibleQuote{罪的\Emph{工资}是死亡}，以表明死亡不是不公正地施加的，而是罪的应得报酬。但恩赐，除非完全是白得的，就不是恩赐。那么我们要明白，人的善功本身就是天主的恩赐，以致当这些善功获得永生的赏报时，那不过是恩宠上加恩宠。因此，人这样被造为正直，以至于虽然缺乏神圣帮助就不能保持正直，却能凭自己的意志离开它。无论他选择哪条道路，天主的旨意都要藉着他或关于他来成就。因此，既然他选择行自己的意愿而非天主的，天主的旨意就在他身上得以完成；因为天主从同一堆毁灭——即人类种族——中，造了一个器皿为尊贵，另一个为卑贱；为尊贵是出于仁慈，为卑贱是出于审判；好叫没有人能因着人而夸耀，因而也不因自己而夸耀。\par

% structure: p0222 source=h3 target=subsection reason=relative-heading-rank; base=h2

\subsection{108、为使天主与我们和好，中保是必需的；除非这中保是天主，他不能是我们的救赎者。}

因为即使通过天主与人之间的唯一中保，那人基督耶稣，我们也不能被救赎，除非他也是天主。当亚当被造时，他是一个义人，并不需要中保。但当罪恶在天主与人类之间造成巨大鸿沟时，便需要一个中保——他唯独在人类中无罪恶地出生、生活且死亡——来使我们与天主和好，并且为我们的身体也获得复活至永生，以便藉着天主的谦卑揭露并治愈人的骄傲；好叫人看出他离天主有多远，当时天主降生成人使他回归；好给不顺从的人树立一个榜样，即天主之人的顺服生活；好让恩宠的泉源由独生子取了奴仆的形体（这形体没有任何先前的功德）而开启；好让所应许给被救赎者的身体复活的保证，在救赎者的复活中赐下；好让魔鬼被它夸耀所欺骗的同一个本性所征服，而人却不被荣耀，免得骄傲再次抬头；总之，为了一切明达者所能感知和描述、或感知而无法描述的种种益处，这些益处源自中保位格的超越奥秘。\par

% structure: p0224 source=h3 target=subsection reason=relative-heading-rank; base=h2

\subsection{109、灵魂在死亡与复活之间的状态。}

此外，在一个人死亡与最终复活之间的期间，灵魂居于一个隐秘的处所，根据它在地上所过生活赚得的功德，享受安息或遭受痛苦。\par

% structure: p0226 source=h3 target=subsection reason=relative-heading-rank; base=h2

\subsection{110、活着的朋友以圣事和施舍为亡者灵魂带来的益处。}

不可否认，亡者的灵魂因着他们活着的亲友的虔诚而获得益处——这些亲友代表他们奉献中保的祭献，或在教会中为他们施舍。但这些善工只对那些在世时已赢得如此功劳的人有益，以致这类善工能帮助他们。因为有一种生活方式，既非好到死后不需要这类善工，也非坏到死后这类善工完全无效；另一方面，有一种生活好到不需要它们；还有一种生活坏到生命结束后它们毫无帮助。因此，所有能减轻或加重人死后痛苦的功德或罪过，都是在此生获得的。所以，没有人可以希望在死后获得他在此世忽略去争取的天主的垂青。由此可见，教会为亡者所举行的善工，绝不与宗徒的话相抵触：\BibleQuote{因为我们众人都必须出现在基督的审判台前，好使每人按他在肉身内所行的，或善或恶，领取相当的报应。}因为使这类善工对人有益的那些功劳，是在他活在肉身内时挣得的。这些善工并非对每个人都有益。为什么它们并非对所有人有益？岂不正是因为人们在肉身内所过的生活种类不同吗？因此，当为所有已受洗的亡者奉献祭台上的祭献或施舍时，它们是对极好的人谢恩的祭献，是对不太坏的人赎罪的祭献；至于极坏的人，纵然它们对亡者无济于事，但对活着的人却是一种安慰。而在那些有益的情况下，其益处在于获得罪过的完全赦免，或至少使惩罚较为可忍。\par

% structure: p0228 source=h3 target=subsection reason=relative-heading-rank; base=h2

\subsection{第111章 复活后，将有两个截然不同的国度：一个是永恒的幸福，另一个是永恒的悲惨。}

然而，在复活之后，当最后普世审判完成后，将有两个国度，各有其明确的边界：一个是基督的，另一个是魔鬼的；一个由善人组成，另一个由恶人组成——但两者都由天使和人组成。前者将没有犯罪的意愿，后者将没有犯罪的能力；两者都无权选择死亡；但前者将在永生中真实而幸福地生活，后者将在永死中悲惨地苟延残喘，却无法死去；因为生命和死亡都将没有终结。但在前者之中，将有不同程度的幸福，一人比另一人更为卓越地幸福；在后者之中，将有不同程度的痛苦，一人比另一人较为可忍受地痛苦。\par

% structure: p0230 source=h3 target=subsection reason=relative-heading-rank; base=h2

\subsection{第112章 圣经中毫无根据来支持那些否认未来惩罚永恒性之人的意见。}

因此，有些人，实际上很多人，哀叹被罚者所受的永恒惩罚和永不止息、无间断的痛苦，并说他们不相信事情会是如此；他们并非直接反对圣经，而是受自身情感的建议，对一切看似严厉的言辞加以缓和，并给予一种更温和的理解，认为这些陈述与其说应照字面接受，不如说旨在使人畏惧。他们说：\BibleQuote{天主岂忘记了施恩？}祂在忿怒中关闭了祂的仁慈吗？他们在一篇圣咏中读到这句话。但无疑，我们应理解这句话是指那些在别处被称为\BibleQuote{蒙受怜悯的器皿}的人说的，因为连他们脱离痛苦也不是凭自己的任何功德，而纯粹是天主的怜悯。或者，如果我们所说的这些人坚持这段经文适用于全人类，就没有理由因此认为那些被说\BibleQuote{这些要进入永罚}的人的惩罚会有终结；因为这将与被说\BibleQuote{义人却要进入永生}的人的幸福以同样方式、在同一时刻终结。但如果他们喜欢这样想，就让他们假设被罚者的痛苦在某些间隔期有所减轻。因为即使在这种情况下，天主的愤怒，即他们的定罪（因为被称为天主愤怒的是这个，而非天主思想中的任何搅动情绪），仍停留在他们身上；也就是说，祂的愤怒虽仍存留，却并非关闭祂的仁慈；虽然祂的仁慈彰显出来，不是要终结他们的永恒惩罚，而是要减轻或给予他们痛苦中的喘息；因为圣咏并没有说\BibleQuote{结束祂的愤怒}或\BibleQuote{当祂的愤怒过去时}，而是说\BibleQuote{在祂的愤怒中}。现在，如果这愤怒是单独的，或者它存在于最小可能程度，那么，失落于天主的国之外，被流放离开天主的城，与天主的生命隔绝，无分于天主为敬畏祂的人所预备、为信赖祂的人所成就的那极大美善，将是一种如此巨大的惩罚，假设它是永恒的，那么我们所知的任何痛苦，即使延续人类想象力所能设想的许多世代，都不能与它相比。\par

% structure: p0232 source=h3 target=subsection reason=relative-heading-rank; base=h2

\subsection{第113章 恶人的死与圣人的生一样将是永恒的。}

因此，恶人这种永久的死亡，即他们与天主的生命隔绝，将永远持续，并且是所有人共有的，无论人们因着人的情感而对各种惩罚或对痛苦的减轻或暂停做出何种猜想；正如圣人的永生将永远持续，并且是所有人共有的，无论那些在和谐光辉中照耀的人们之间有多少等级和荣誉。\par

% structure: p0234 source=h2 target=omitted reason=duplicate-page-title

% structure: p0235 source=h3 target=subsection reason=relative-heading-rank; base=h2

\subsection{第114章 谈论信德后，我们现在来讨论望德。凡属于望德的一切，都包含在主祷文中。}

由这信德的宣认（它被简要地概括在信经中，按血肉理解是给婴孩的奶，但按精神领悟和研读则是给强壮者的干粮），产生了信徒的善\Emph{望德}；而这望德伴随着神圣的\Emph{爱德}。但在这些全是真实信德对象的事中，只有那些包含在主祷文中的才涉及望德。因为\Quote{信赖世人的人，是可咒骂的}是圣经的见证；因此，这咒骂也归于信赖自己的人。所以，除非从主天主那里，我们不应祈求任何我们盼望善行或盼望作为善工赏报得到的事物。\par

% structure: p0237 source=h3 target=subsection reason=relative-heading-rank; base=h2

\subsection{第115章 主祷文的七个祈求，按玛窦福音。}

因此，在玛窦福音中，主祷文似乎包含七个祈求，其中三个求永恒福乐，其余四个求现世福乐；然而后者是达到永恒福乐所必需的先决条件。因为当我们说\Quote{愿你的名被尊为圣；愿你的国来临；愿你的旨意奉行在人间，如同在天上}（有人并不无理地解释为在身体和灵魂中），我们祈求的是永远享受的福乐；这些福乐确实在此世开始，并随着我们在恩宠中成长而增长，但在其完美状态（当在来生中期待）中，将成为永久的产业。但当我们说\Quote{求你今天赏给我们日用的食粮；并宽免我们的罪债，犹如我们也宽免得罪我们的人；不要让我们陷于诱惑，但救我们免于凶恶}，谁看不出我们祈求的是与今生需要有关的福乐呢？在永生中（我们盼望永远生活在那里），天主之名的圣化、祂的国和祂在我们灵魂与身体中的旨意，将达致完美，并持续到永远。而我们的\Emph{日用}食粮之所以如此称呼，是因为在此常有对精神与肉体所需营养的需求，无论我们按精神、按肉体或按双重意义来理解。也正是在此处，我们需要祈求宽恕，因为我们在此犯罪；此处有诱惑，诱使我们或驱使我们犯罪；简言之，此处有我们渴望从中解脱的凶恶：但在另一个世界，这些都不存在。\par

% structure: p0239 source=h3 target=subsection reason=relative-heading-rank; base=h2

\subsection{第116章 路加以更简洁的方式将这七个祈求的实质表达为五个。}

但圣史路加在他的主祷文版本中并未包含七个，而是五个祈求：当然，这不是说两位圣史之间有矛盾，而是路加以他的简洁表明了玛窦的七个祈求应如何理解。因为天主的名在精神中被尊为圣；天主的国将在身体复活时来临。因此，路加意图表明第三个祈求是对前两个的一种重复，他选择完全省略第三个祈求来指明这一点。然后他添加了另外三个：一个为日用食粮，另一个为罪的赦免，另一个为免受诱惑。而玛窦列为最后一个祈求的\Quote{但救我们免于凶恶}，路加省略了，以向我们表明它已包含在前一个关于诱惑的祈求中。实际上，玛窦自己说\Quote{\Emph{但}救我们}，而不是\Quote{\Emph{并}救我们}，仿佛要表明这两个祈求实际上是一个：不要做这个，而要这个；以便每个人都能理解，在他不被引入诱惑这一事实中，他就被救离凶恶。\par

% structure: p0241 source=h2 target=omitted reason=duplicate-page-title

% structure: p0242 source=h3 target=subsection reason=relative-heading-rank; base=h2

\subsection{第117章 爱德，比信德和望德更大，由圣神倾注在我们心中。}

现在论及\Emph{爱德}，宗徒宣称它大于其他两种恩赐，即信德和望德；爱德越丰富地居住在一个人内，这人也越美好。因为当问及一个人是否良善时，人们不问他的信仰或望德，而问他爱什么。因为正确去爱的人无疑也正确相信和盼望；而没有爱的人即使他的信仰真实，也是徒然相信；即使他盼望的对象真正属于真福的一部分，也是徒然盼望；除非他相信并盼望藉祈祷得到爱德的恩赐。因为虽然不可能没有爱而盼望，但有可能一个人不爱那达到他盼望所必需的事物；例如，如果他盼望永生（谁不渴望这个呢？）却又不爱正义，没有正义无人能达至永生。这就是宗徒所说的基督的真信德\Quote{藉爱德行事的}；如果它在其爱中尚未拥有任何东西，它就祈求领受，寻求找到，敲门而得以开启。因为信德藉祈祷获得法律所命令的。因为没有天主的恩赐，即没有圣神（我们心中的爱是由祂倾注的），法律能够命令，却不能帮助；而且，它使人成为违犯者，因为他不能再以无知为借口。如今，哪里没有天主的爱，肉欲就在那里作王。\par

% structure: p0244 source=h3 target=subsection reason=relative-heading-rank; base=h2

\subsection{第118章 基督徒生活的四个阶段，以及教会历史相对应的四个阶段。}

当人沉溺在无知的至深黑暗之中，按血肉生活，不受理性或良心的任何挣扎所扰时，这便是他的第一种状态。随后，借着法律，对罪的认识来到，而天主之神尚未施予祂的援助，人努力按法律生活，却受阻而失败，陷入明知之罪，于是被罪制伏，成了它的奴隶（\Quote{因为人被谁制伏，就是谁的奴仆}\BibleRef{伯后2:19}）；由此，对诫命的认识所产生的效果是：罪在人内激起各种贪欲，他更犯上故意违命的罪债，这就应验了经上所记载的：\BibleQuote{法律本是外添的，为叫过犯显多}\BibleRef{罗5:20}。这是人的第二种状态。但如果天主眷顾他，感动他信赖天主的帮助，天主之神开始在他内工作，那么爱德的更大力量便与血肉之力相争；虽然人的本性中仍有一种与他交战的力量（因为他的病患尚未完全痊愈），但他因信德而活义人生活，并活于正义，只要他不屈服于邪恶的贪欲，而因对圣洁的爱战胜它。这是怀有良好望德之人的第三种状态。凡以坚定的虔敬在这过程中前进的人，终将获得平安，这平安在今生结束后，将在灵魂的安息中得以成全，最终在身体的复活中达至圆满。在这四种不同的阶段中，第一种是在法律之前，第二种是在法律之下，第三种是在恩宠之下，第四种是在完全而圆满的平安中。同样，天主子民的历史也按祂的旨意安排，祂按数目、分量和轻重安排万有。因为教会最初存在于法律之前；然后在法律之下，这法律是由\Person{梅瑟}颁布的；然后在恩宠之下，恩宠首先在\OriginalTerm{中保}{Mediator}来临时显明。当然，这恩宠并非先前不存在，而是按时代的安排，被遮掩隐藏了。因为没有一个人，即使是古代义人，能离了基督的信德而得救恩；除非他们认识祂，他们的职务也不能被用来向我们传递有关祂的预言，有些更明显，有些更隐晦。\par

% structure: p0246 source=h3 target=subsection reason=relative-heading-rank; base=h2

\subsection{第119章　\OriginalTerm{重生之恩}{The Grace of Regeneration}洗去一切过去的罪和一切\OriginalTerm{原罪之债}{Original Guilt}}

现在，无论这四种阶段（我们可这样称呼）中的哪一种，\OriginalTerm{重生}{regeneration}之恩在遇见任何特定的人时，他所有过去的罪都立即被赦免，他在出生时所得的\OriginalTerm{罪债}{guilt}也在新生中被除去；\BibleQuote{风随意向哪里吹}\BibleRef{若3:8}，这话千真万确，以致有些人从未经历过第二种阶段，即法律下的奴役，而是一接受诫命，就领受了天主的援助。\par

% structure: p0248 source=h3 target=subsection reason=relative-heading-rank; base=h2

\subsection{第120章　死亡不能伤害那些已领受\OriginalTerm{重生之恩}{The Grace of Regeneration}的人}

但在人能领受诫命之前，他必须按血肉生活。然而，一旦他领受了\OriginalTerm{重生}{regeneration}之恩，死亡就不能伤害他，即使他立刻离开此生；\BibleQuote{因为基督死而复活，正是为作死人和活人的主}\BibleRef{罗14:9}；死亡也不能辖制那基督已为他白白死去的人。\par

% structure: p0250 source=h3 target=subsection reason=relative-heading-rank; base=h2

\subsection{第121章　爱德是一切诫命的终向，天主自己就是爱}

天主的一切诫命都包含在爱德之中，宗徒说：\BibleQuote{诫命的终向是爱德，出于纯洁的心、良好的良心和真诚的信德}\BibleRef{弟前1:5}。因此，每条诫命的终向都是爱德，也就是说，每条诫命都以爱为其目标。但凡是出于对惩罚的恐惧或某种其他属肉体的动机而行的，其原则不是天主之神倾注在心里的爱，那么无论在人看来如何，都不算按应当的方式而行。因为这爱德包含对天主的爱和对近人的爱，\BibleQuote{全部法律和先知都系于这两条诫命}\BibleRef{玛22:40}，我们还可以加上福音和宗徒。因为我们正是从这些听到这声音：诫命的终向是爱德，天主就是爱。因此，天主的一切诫命，其中一条是\BibleQuote{不可奸淫}\BibleRef{出20:14}，以及所有那些不是诫命而是特殊劝谕的规条，其中一条是\BibleQuote{男人不接触女人是好的}\BibleRef{格前7:1}，只有行动的动机原则是对天主的爱以及对近人并在天主内的爱时，才能正确实行。这适用于今生和来生。我们现在因信德爱天主，那时我们将借着看见而爱祂。我们现在甚至因信德爱我们的近人；因为我们是会死的人，不知道必死之人的心。但在来生，主\BibleQuote{要暴露暗中的隐情，并揭示人心的计谋，那时各人才从天主那里获得称赞}\BibleRef{格前4:5}；因为每人要在邻人身上爱和赞美那美德，这美德为了不被隐藏，主要亲自彰显。此外，\OriginalTerm{贪欲}{concupiscence}随着爱的增长而减少，直到后者增长到在此生不能再增长的高度。因为\BibleQuote{人若为自己的朋友舍命，没有比这更大的爱}\BibleRef{若15:13}。谁能说将来世界上的爱要有多大呢？那时不再有\OriginalTerm{贪欲}{lust}需要它克制和战胜？因为那将是健康的完美，不再有与死亡的斗争。\par

% structure: p0252 source=h3 target=subsection reason=relative-heading-rank; base=h2

\subsection{第122章　结论}

但现在，这卷书终于要结束了。由您自己判断是否可以称它为一本手册，或将其用作手册。而我，因着我认为不应轻视您对基督的热忱，并依赖我们\OriginalTerm{救主}{Redeemer}的帮助，对您怀有并盼望一切美善，且深爱您作为祂身体的一个肢体，我尽我所能为您写了这本关于\WorkTitle{信德、望德与爱德}的书。愿它的价值与其篇幅相当。\par

\SourceRecord{Translated by J.F. Shaw.；Nicene and Post-Nicene Fathers, First Series；Vol. 3.；Edited by Philip Schaff.；Buffalo, NY: Christian Literature Publishing Co.,；1887.；<http://www.newadvent.org/fathers/1302.htm>.}
